% =============================================================================
% MYCELLIUM-AEGIS · SAYISAL FİZİBİLİTE ANALİZİ
% Derleme:  pdflatex -shell-escape yok, iki geçiş yeterli
%   1) python numeric/run_all.py      (data/*.dat + summary.json)
%   2) python numeric/geo.py          (data/geo_*.dat + summary_geo.json)
%   3) python numeric/make_tikz.py    (figures/*.tikz)
%   4) python numeric/make_numbers.py (numbers.tex)
%   5) pdflatex ... ×2
% =============================================================================
\documentclass[11pt,a4paper]{article}

\usepackage[utf8]{inputenc}
\usepackage[T1]{fontenc}
\usepackage[shorthands=off,turkish]{babel}
\usepackage{lmodern}
\usepackage[a4paper,margin=2.4cm]{geometry}
\usepackage{amsmath,amssymb,amsthm,bm}
\usepackage{booktabs}
\usepackage{array}
\usepackage{longtable}
\usepackage{graphicx}
\usepackage{caption}
\usepackage{microtype}
\usepackage[hidelinks]{hyperref}

% =============================================================================
% MYCELLIUM-AEGIS · FİZİBİLİTE ANALİZİ
% preamble.tex — pgfplots ortak biçemi
% Hem ana rapor hem de bağımsız (standalone) şekil derlemeleri bunu yükler.
% =============================================================================
\usepackage{pgfplots}
\usepackage{pgfplotstable}
\usepgfplotslibrary{groupplots,fillbetween,colorbrewer,statistics}
\usetikzlibrary{patterns,arrows.meta,calc,decorations.pathreplacing}
\pgfplotsset{compat=1.18}

% ---- proje renk dizgesi (sunumla aynı) --------------------------------------
\definecolor{aegisteal}{HTML}{0E8F76}
\definecolor{aegisdeep}{HTML}{0A6B58}
\definecolor{aegisember}{HTML}{C1670E}
\definecolor{aegisink}{HTML}{1E2622}
\definecolor{aegismuted}{HTML}{6B7A72}
\definecolor{aegisline}{HTML}{D6DBD6}
\definecolor{aegispaper}{HTML}{FBFBF9}
\definecolor{aegisblue}{HTML}{2C5F8A}
\definecolor{aegisplum}{HTML}{7A4A78}

% arazi yüksekliği için sessiz, projeye uyumlu bir renk skalası
\pgfplotsset{
  colormap={aegisterrain}{
    rgb255=(232,238,232) rgb255=(200,214,198) rgb255=(178,190,160)
    rgb255=(196,182,140) rgb255=(190,158,116) rgb255=(168,128,96)
    rgb255=(142,104,84)  rgb255=(120,88,74)
  },
}

\pgfplotscreateplotcyclelist{aegis}{%
  {aegisteal, thick},
  {aegisember, thick, dashed},
  {aegisblue, thick, dotted},
  {aegisplum, thick, dashdotted},
  {aegismuted, thick, densely dashed},
}

% ---- taban eksen biçemi -----------------------------------------------------
\pgfplotsset{
  aegis/.style={
    width=\figwidth, height=\figheight,
    scale only axis,
    axis line style={aegisink, line width=.5pt},
    tick style={aegisink, line width=.4pt},
    label style={font=\small, color=aegisink},
    tick label style={font=\footnotesize, color=aegisink},
    title style={font=\small\bfseries, color=aegisink},
    legend style={font=\footnotesize, draw=aegisline, fill=aegispaper,
                  fill opacity=.92, text opacity=1, rounded corners=1pt,
                  inner sep=3pt, row sep=1pt},
    grid=major,
    grid style={aegisline, line width=.3pt},
    cycle list name=aegis,
    every axis plot/.append style={line join=round},
    clip mode=individual,
  },
  aegismap/.style={
    aegis, grid=none, axis equal image,
    xlabel={Doğu [km]}, ylabel={Kuzey [km]},
    scaled ticks=false,
    xticklabel style={/pgf/number format/.cd, fixed, precision=0, 1000 sep={\,}},
    yticklabel style={/pgf/number format/.cd, fixed, precision=0, 1000 sep={\,}},
    colorbar style={font=\scriptsize, width=6pt, tick style={line width=.3pt}},
  },
  /pgf/number format/.cd, use comma, 1000 sep={\,},
}

% Türkçe ondalık ayırıcı için sayı biçimi
\pgfkeys{/pgf/number format/precision=3}

\newcommand{\figwidth}{12.2cm}
\newcommand{\figheight}{6.4cm}

% veri kökü — ana rapor ve standalone derlemede aynı olsun diye
\providecommand{\dataroot}{data}

% OTOMATİK ÜRETİLDİ — numeric/make_numbers.py · elle düzenlemeyin
\newcommand{\tKpar}{1,60}
\newcommand{\tKperp}{0,85}
\newcommand{\tAniso}{1,88}
\newcommand{\tMisMax}{17,8}
\newcommand{\tDipMax}{36}
\newcommand{\tKxzTf}{-0,287}
\newcommand{\pThermal}{2,03}
\newcommand{\pErt}{2,00}
\newcommand{\pPoint}{1,01}
\newcommand{\pGrid}{0,88}
\newcommand{\gciGrid}{0,74}
\newcommand{\pTime}{1,01}
\newcommand{\dtExplicit}{2,2}
\newcommand{\coercivity}{3{,}16\times 10^{-8}}
\newcommand{\dblkOO}{1{,}45\times 10^{-6}}
\newcommand{\dblkOT}{7{,}85\times 10^{-10}}
\newcommand{\dblkTO}{4{,}87\times 10^{-10}}
\newcommand{\dblkTT}{3{,}16\times 10^{-8}}
\newcommand{\stiffRatio}{46}
\newcommand{\latentOff}{79,2}
\newcommand{\latentOn}{79,2}
\newcommand{\latentDelta}{0,0}
\newcommand{\leadRef}{79,7}
\newcommand{\numUnc}{0,4}
\newcommand{\leadCoarse}{77,5}
\newcommand{\leadFine}{79,7}
\newcommand{\regSpread}{0,0}
\newcommand{\regLo}{79,3}
\newcommand{\regHi}{79,3}
\newcommand{\leadNom}{79,3}
\newcommand{\tAlarm}{3,18}
\newcommand{\tIgnite}{4,50}
\newcommand{\maxdTdt}{23,8}
\newcommand{\mindlnR}{0,0419}
\newcommand{\dampDepth}{15,6}
\newcommand{\diurnalRate}{0,36}
\newcommand{\betaMarginNat}{4,2}
\newcommand{\droughtdT}{0,55}
\newcommand{\droughtdlnR}{0,0012}
\newcommand{\normaldT}{0,35}
\newcommand{\normaldlnR}{0,0009}
\newcommand{\sepBeta}{43}
\newcommand{\leadSix}{79,3}
\newcommand{\leadEighteen}{24,9}
\newcommand{\zZero}{21,6}
\newcommand{\zOneMax}{5}
\newcommand{\leadOneBest}{22,1}
\newcommand{\zOneBest}{2}
\newcommand{\zMedian}{6,0}
\newcommand{\zPten}{2,6}
\newcommand{\zPninety}{9,4}
\newcommand{\rPninety}{4,1}
\newcommand{\bornErr}{3,0}
\newcommand{\mcPos}{1024}
\newcommand{\mcNeg}{1024}
\newcommand{\detRate}{100,0}
\newcommand{\leadPfive}{24,9}
\newcommand{\leadMed}{76,6}
\newcommand{\leadPnf}{157,7}
\newcommand{\pLeadFifteen}{99,2}
\newcommand{\pLeadThirty}{92,2}
\newcommand{\aucFusion}{1,000}
\newcommand{\aucBeta}{1,000}
\newcommand{\aucGamma}{1,000}
\newcommand{\tprOne}{100,0}
\newcommand{\fprOne}{0,00}
\newcommand{\betaLo}{0,61}
\newcommand{\betaHi}{5,38}
\newcommand{\gammaLo}{0,0013}
\newcommand{\gammaHi}{0,0098}
\newcommand{\sobolRuns}{5120}
\newcommand{\sobolN}{512}
\newcommand{\sobolSum}{0,98}
\newcommand{\sobolFirst}{piroliz ramp süresi}
\newcommand{\sobolFirstST}{0,48}
\newcommand{\sobolSecond}{yüzey plato sıcaklığı}
\newcommand{\sobolSecondST}{0,43}
\newcommand{\sobolThird}{prob derinliği}
\newcommand{\sobolThirdST}{0,09}
\newcommand{\sobolNull}{0,000}
\newcommand{\latSpot}{0,60}
\newcommand{\latRzero}{0,67}
\newcommand{\latRup}{0,69}
\newcommand{\latRdown}{0,65}
\newcommand{\cableVzero}{30}
\newcommand{\cableVdet}{0,5}
\newcommand{\rdetHalf}{1,3}
\newcommand{\rdetOne}{2,4}
\newcommand{\rdetThree}{6,7}
\newcommand{\rdetTen}{20,6}
\newcommand{\rdetThirty}{58,0}
\newcommand{\covNodes}{1,15}
\newcommand{\covCapex}{5 669}
\newcommand{\covArea}{0,82}
\newcommand{\covTime}{43}
\newcommand{\nodeTL}{4 910}
\newcommand{\iavgFifty}{4,77}
\newcommand{\iavgTen}{1,06}
\newcommand{\lolpFifty}{0,00}
\newcommand{\lolpTen}{0,00}
\newcommand{\lolpShade}{0,0}
\newcommand{\canopyCritFifty}{3,2}
\newcommand{\canopyCritFull}{7,7}
\newcommand{\energyMargin}{11,8}
\newcommand{\pshWinter}{1,74}
\newcommand{\pshSummer}{7,50}
\newcommand{\dMax}{2 639}
\newcommand{\dMarginTen}{1 285}
\newcommand{\nodesPerGw}{599}
\newcommand{\nodesRange}{599}
\newcommand{\nodesAirtime}{9 600}
\newcommand{\pathExp}{3,2}
\newcommand{\npvMed}{6,4}
\newcommand{\npvPten}{-4,8}
\newcommand{\npvPninety}{27,3}
\newcommand{\pNpvPos}{73,7}
\newcommand{\irrMed}{62}
\newcommand{\irrPten}{21}
\newcommand{\paybackYear}{6}
\newcommand{\econRuns}{20 000}
\newcommand{\geoLandArea}{13 900}
\newcommand{\geoElevMax}{2 136}
\newcommand{\geoSlopeMean}{8,1}
\newcommand{\geoHighArea}{1 338}
\newcommand{\geoRes}{90}
\newcommand{\geoCells}{4000}
\newcommand{\geoLeadMed}{78,1}
\newcommand{\geoLeadPfive}{74,6}
\newcommand{\geoLeadPnf}{79,8}
\newcommand{\geoTopProv}{Aydın}
\newcommand{\geoTopRisk}{0,639}
\newcommand{\geoSecondProv}{İzmir}
\newcommand{\geoNprov}{3}
\newcommand{\geoTransLen}{127}
\newcommand{\geoTransZmax}{1 165}
\newcommand{\geoTransLead}{78,4}
\newcommand{\geoSiteN}{8}
\newcommand{\geoSiteRisk}{0,889}
\newcommand{\geoSiteElev}{1 077}
\newcommand{\geoSiteSlope}{27,0}


\captionsetup{font=small, labelfont={sc}, width=.92\linewidth}
\setlength{\parskip}{0.35em}
\renewcommand{\arraystretch}{1.15}

\newcommand{\dd}{\mathrm{d}}
\newcommand{\Div}{\nabla\!\cdot\!}
\newcommand{\bK}{\mathbf{K}}
\newcommand{\bq}{\mathbf{q}}
\newcommand{\bn}{\mathbf{n}}
\newcommand{\bu}{\mathbf{u}}
\newcommand{\bD}{\bm{\mathcal{D}}}
\newcommand{\bsig}{\bm{\sigma}}

\title{\vspace{-1.5cm}\Large Mycellium-Aegis\\[.25em]
\huge Sayısal Fizibilite Analizi\\[.5em]
\large\normalfont Bağlaşık tensör--kısmi diferansiyel denklem modeli,\\
doğrulanmış sayısal çözüm ve coğrafi uygulanabilirlik}
\author{\normalsize TÜBİTAK 1812 BİGG · AGY-112 eki}
\date{\normalsize Temmuz 2026}

\begin{document}
\maketitle
\thispagestyle{empty}

\begin{abstract}\noindent
Bu belge, Mycellium-Aegis erken uyarı sisteminin yapılabilirliğini niteliksel
gerekçelerle değil, çözülmüş bir fiziksel modelin sayısal sonuçlarıyla
değerlendirir. Toprakta ısı ve nemin bağlaşık taşınımı, ikinci mertebe bir
iletkenlik tensörü ve dört indisli bir yayınım nesnesi üzerinden kurulan bir
kısmi diferansiyel denklem sistemi olarak yazılmış; sonlu hacim ayrıklaştırması
üretilmiş çözümlerle doğrulanmış; ölçülen elektriksel büyüklük, eksenel simetrik
bir potansiyel probleminin Fréchet türevi üzerinden tanımlanmıştır. Sistemin
alarm kuralı, \mcPos\ örneklik bir Monte Carlo tasarımı ve \sobolRuns\ koşumluk
bir Sobol duyarlılık çözümlemesiyle sınanmış; enerji, haberleşme ve ekonomik
katmanlar aynı belirsizlik çerçevesine bağlanmıştır. Son bölüm, modeli İzmir
Orman Bölge Müdürlüğü kapsama alanının otuz metrelik sayısal yükseklik modeli
üzerine oturtarak öncü sürenin mekânsal dağılımını ve kesitler üzerinden
yeraltı sıcaklık alanını verir. Analizin iki temel bulgusu şudur: belgelerde
tanımlı tek yönlü direnç kuralı, prob derinleştikçe öncü sürenin neredeyse
tamamını yitirmektedir ve mutlak değer üzerinden yeniden yazılmalıdır; sistemin
mekânsal kapsama fizibilitesi ise miselyum ağının elektrotonik uzunluk sabitine
doğrudan bağlıdır ve bu büyüklüğün ölçülmesi WP-2'nin birincil çıktısı
olmalıdır.
\end{abstract}

\vspace{.4em}
{\small\tableofcontents}
\newpage

% =============================================================================
\section{Kapsam, soru ve yöntem}

Bir erken uyarı sisteminin fizibilitesi tek bir soru değildir. Sistemin
ölçmeyi iddia ettiği fiziksel büyüklüğün, alev görünmeden önce ölçülebilir
bir değişim gösterip göstermediği bir sorudur; bu değişimin, aynı toprakta
yangın olmadan da meydana gelen değişimlerden ayrılabilir olup olmadığı ikinci
ve bağımsız bir sorudur. Ölçüm zincirinin bu ayrımı taşıyacak duyarlılığa ve
enerjiye sahip olup olmadığı üçüncü, kurulacak düğüm yoğunluğunun kapsamak
istenen alanı ekonomik biçimde örtüp örtmediği dördüncü sorudur. Bu belge
dördünü de aynı fiziksel modelden türetilen sayılarla ele alır ve her birine
sayısal bir eşik atar.

Yöntem tercihinin gerekçesi şudur. Toprak altında alev öncesi ısınmanın ve
kurumanın birbirine bağlı iki taşınım süreci olduğu, laboratuvar ölçeğinde
zaten gözlenmiştir; ancak bu iki sürecin bir prob derinliğinde hangi sırayla
ve hangi işaretle görüneceği, kurucu bağıntıların birlikte çözülmesini
gerektirir. Fenomenolojik bir kuruma katsayısı, ölçülen büyüklüğün işaretini
baştan varsaydığı için, tam da fizibilitenin dayandığı noktayı ispat yerine
kabul eder. Bu yüzden burada, buhar taşınımını ve gizli ısıyı içeren bağlaşık
sistemin kendisi çözülmüş, kurucu bağıntılar literatürden alınmış, çözücü
üretilmiş çözümlerle doğrulanmış ve sonuçların ayrıklaştırma seçimlerine
duyarsızlığı ölçülmüştür.

Bütün sayısal çalışma \texttt{docs/feasibility/numeric/} altındadır ve sabit
tohumludur. Metindeki her sayı, \texttt{summary.json} dosyasından üretilen
makrolarla yerleştirilmiştir; çözücü yeniden koşturulduğunda metin de
güncellenir. Şekiller pgfplots ile aynı veri dosyalarından çizilir, yani
gövde metni ile grafikler arasında elle kopyalanmış hiçbir sayı yoktur.

% =============================================================================
\section{Fiziksel model: tensör formülasyonu}

\subsection{Bağlaşık ısı ve nem taşınımı}

Toprak sütununda durum değişkeni, sıcaklık ve hacimsel su içeriğinden oluşan
iki bileşenli bir alan olarak alınır: $\bu = (T,\theta)^{\!\top}$. Philip ve
de Vries'in bağlaşık taşınım kuramında bu iki alan, ayrı ayrı difüzyon
denklemleri değil, çapraz terimlerle bağlanmış bir sistem sağlar:
\begin{equation}
\frac{\partial u_a}{\partial t}
   \;=\; \partial_i\!\left( \mathcal{D}_{abij}\,\partial_j u_b \right) + s_a ,
\qquad a,b\in\{1,2\},\quad i,j\in\{1,2,3\},
\label{eq:coupled}
\end{equation}
burada toplama yinelenen indisler üzerinden yapılır. Dört indisli
$\mathcal{D}_{abij}$ nesnesi, alanlar arası bağlaşımı taşıyan $2\times2$ bir
blok ile uzaysal anizotropiyi taşıyan $3\times3$ bir tensörün çarpımına ayrılır:
\begin{equation}
\mathcal{D}_{abij} \;=\; D_{ab}\,\hat{K}_{ij},
\qquad
D \;=\;
\begin{pmatrix}
\dfrac{\lambda_{\mathrm{ef}}(\theta,T)}{C(\theta)} &
\dfrac{L\rho_w D_{\theta v}(\theta,T)}{C(\theta)}\\[1.1em]
D_{Tv}(\theta,T) & D_{\theta \ell}(\theta) + D_{\theta v}(\theta,T)
\end{pmatrix}.
\label{eq:Dblock}
\end{equation}
Köşegen dışı terimler modelin özüdür. Sağ üstteki terim, nem gradyanının
buhar taşınımı yoluyla taşıdığı gizli ısıyı; sol alttaki terim ise sıcaklık
gradyanının suyu sürüklemesini temsil eder. İkincisi işaret bakımından
belirleyicidir: buhar sıcaktan soğuğa göç ettiği için, yüzeyde ısınan bir
toprakta ısıl cephenin \emph{önünde} bir yoğuşma bölgesi oluşur ve prob
derinliğindeki su içeriği, kuruma cephesi oraya varmadan önce artar. Bu
davranış dördüncü bölümdeki en önemli bulgunun kaynağıdır.

Kurucu bağıntılar literatürden alınmıştır. Su tutma eğrisi ve doymamış
hidrolik iletkenlik van Genuchten--Mualem biçiminde, parametreleri Carsel ve
Parrish'in kumlu tın sınıfına aittir; ısıl iletkenlik Côté ve Konrad'ın
normalize modeliyle, $\lambda(\theta) = \lambda_{\mathrm{kuru}} +
(\lambda_{\mathrm{doygun}}-\lambda_{\mathrm{kuru}})\,\kappa S_r/(1+(\kappa-1)S_r)$
olarak; buhar yayınımları de Vries'in $D_{Tv}$ ve $D_{\theta v}$ ifadeleriyle,
Millington--Quirk tortuozitesi ve Cass ve arkadaşlarının artırma katsayısıyla
hesaplanır. Gizli ısı taşınımı, ısı denkleminde etkin iletkenliği
$\lambda_{\mathrm{ef}} = \lambda + L\rho_w D_{Tv}$ biçiminde büyütür; sıcak ve
nemli toprakta bu terim iletimden büyük olabilir ve modelin ısıyı derine
taşıyan asıl mekanizmasıdır.

Yüz santigrat derecede sıvı suyun ortamda kalamayacağı, ayrı bir kaynama
kinetiğiyle zorlanır. Birim hacimde su kaybı hızı
$\dot S = k_0\,\max(0,T-T_b)\,(\theta-\theta_{\mathrm{kuru}})/\Delta T$ alınır
ve ısı denkleminden $L\rho_w \dot S$ kadar gizli ısı çekilir. Bu, ölçümlerde
belgelenmiş yüz derece civarındaki sıcaklık düzlüğünü üretir; düzlüğün süresi
$k_0$'a değil enerji arzına, yani $L\rho_w \Delta\theta / q$ oranına bağlıdır.
$k_0$ ve $\Delta T$ sayısal gevşeme parametreleridir ve sonuçların bunlara
duyarsızlığı Bölüm~\ref{sec:reg}'te ölçülmüştür.

\subsection{Anizotropi, yatak eğimi ve akı yönelim sapması}

Orman toprağı yatay katmanlıdır; ısıl iletkenlik katman düzleminde ve düzleme
dik yönde farklıdır. Bu, yatak normali $\bn$ etrafında enine izotrop bir
ikinci mertebe tensörle temsil edilir:
\begin{equation}
\bK(\bn) \;=\; k_\parallel\left(\mathbf{I}-\bn\otimes\bn\right)
              \;+\; k_\perp\,\bn\otimes\bn .
\label{eq:Ktensor}
\end{equation}
Arazi eğimli olduğunda yatak normali de eğilir. Eğim açısı $\psi$ için
$\bn(\psi) = \mathbf{R}(\psi)\,\mathbf{e}_3$ alınırsa tensör
$\bK(\psi) = \mathbf{R}\bK_0\mathbf{R}^{\!\top}$ olarak dönüşür ve bileşenleri
\begin{equation}
K_{xx} = k_\parallel\cos^2\!\psi + k_\perp\sin^2\!\psi,\quad
K_{zz} = k_\parallel\sin^2\!\psi + k_\perp\cos^2\!\psi,\quad
K_{xz} = (k_\perp-k_\parallel)\sin\psi\cos\psi
\end{equation}
biçimini alır. Değişmezler $\operatorname{tr}\bK = k_\parallel+k_\perp$ ve
$\det\bK = k_\parallel k_\perp$ eğimden bağımsızdır; sayısal çözümde bu
özdeğer değişmezliği bir tutarlılık denetimi olarak kullanılmıştır.

Çapraz bileşenin varlığı, ısı akısının sıcaklık gradyanına paralel olmadığı
anlamına gelir. Düşey bir gradyan altında akı ile gradyan arasındaki açı
$\tan\vartheta = |K_{xz}|/K_{zz}$ olup, eğime göre en büyük değerini
\begin{equation}
\tan\vartheta_{\max} = \frac{k_\parallel-k_\perp}{2\sqrt{k_\parallel k_\perp}},
\qquad
\psi^\star = \arctan\sqrt{k_\perp/k_\parallel}
\end{equation}
noktasında alır. $k_\parallel = \tKpar$ ve $k_\perp = \tKperp$
W\,m$^{-1}$K$^{-1}$ için bu, $\psi^\star \approx \tDipMax^\circ$ eğimde
$\vartheta_{\max} = \tMisMax^\circ$ demektir. Şekil~\ref{fig:tensor} bu
davranışı sayısal olarak verir ve analitik en büyük değerle örtüşür.
Uygulamadaki karşılığı somuttur: eğimli arazide bir sıcak nokta, düğümü
düşeyde değil, yamaç yönünde kaydırılmış bir noktada en erken uyarır; bu da
Bölüm~\ref{sec:lateral}'te ölçülen tespit ayak izinin asimetrisini açıklar.

\begin{figure}[htbp]\centering
\begin{tikzpicture}
\begin{axis}[aegis, xlabel={Yatak eğimi $\psi$ [derece]},
  ylabel={Tensör bileşeni [W\,m$^{-1}$K$^{-1}$]},
  legend pos=south west, legend columns=3, xmin=0, xmax=45, ymin=-0.55, ymax=1.85]
\addplot[aegisteal, thick] table[x=dip,y=Kxx] {\dataroot/tensor_dip.dat};
\addlegendentry{$K_{xx}$}
\addplot[aegisblue, thick, dashed] table[x=dip,y=Kzz] {\dataroot/tensor_dip.dat};
\addlegendentry{$K_{zz}$}
\addplot[aegisplum, thick, densely dotted] table[x=dip,y=Kxz] {\dataroot/tensor_dip.dat};
\addlegendentry{$K_{xz}$}
\addplot[aegisline, thin, forget plot] coordinates {(0,0) (45,0)};
\end{axis}
\begin{axis}[aegis, grid=none, axis y line*=right, axis x line=none,
  xmin=0, xmax=45, ymin=0, ymax=26,
  ylabel={$-\nabla T$ ile $\mathbf{q}$ arasındaki açı [derece]},
  ylabel style={aegisember}, y tick label style={aegisember}]
\addplot[aegisember, very thick] table[x=dip,y=misalign] {\dataroot/tensor_dip.dat};
\end{axis}
\end{tikzpicture}

\caption{Enine izotrop ısıl iletkenlik tensörünün bileşenleri ve düşey
sıcaklık gradyanı altında ısı akısının gradyanla yaptığı açı, yatak eğimine
göre. Çapraz bileşen $K_{xz}$ eğimle birlikte doğar ve akıyı yamaç yönünde
sapmaya zorlar.}
\label{fig:tensor}
\end{figure}

\subsection{Elektriksel iletkenlik tensörü ve ölçülen büyüklük}

Sistemin ikinci ölçüm kanalı, elektrot çifti arasındaki transfer direncidir.
Toprağın elektriksel iletkenliği Archie yasasıyla su içeriğine, Arps
düzeltmesiyle sıcaklığa bağlanır:
\begin{equation}
\sigma(\theta,T) \;=\; \sigma_w^{25}\left[1+\alpha_T (T-25)\right]\,
\phi^{m}\left(\frac{\theta}{\phi}\right)^{\! n_A}.
\label{eq:archie}
\end{equation}
Katmanlı ortamda iletkenlik de tensöreldir; düşey simetri ekseni için
$\bsig = \operatorname{diag}(\sigma_h,\sigma_h,\sigma_v)$ ve anizotropi oranı
$\Lambda^2 = \sigma_h/\sigma_v$ yazılır. Elektrot çiftinin gördüğü potansiyel,
\begin{equation}
\Div\!\left(\bsig\nabla\varphi\right) \;=\; -I\left[\delta(\mathbf{x}-\mathbf{x}_+)
- \delta(\mathbf{x}-\mathbf{x}_-)\right]
\label{eq:poisson}
\end{equation}
denkleminin çözümüdür. Bütün elektrotlar tek bir düşey çubuk üzerinde
olduğundan problem eksenel simetriktir ve silindirik koordinatlarda
\begin{equation}
\frac{1}{r}\,\partial_r\!\left(r\,\sigma_h\,\partial_r\varphi\right)
+ \partial_z\!\left(\sigma_v\,\partial_z\varphi\right)
= -\frac{I\,\delta}{2\pi r}
\end{equation}
biçimine iner. Homojen anizotropik yarı-uzayda, $x'=x/\sqrt{\sigma_h}$,
$z'=z/\sqrt{\sigma_v}$ ölçeklemesi denklemi Laplace denklemine indirger ve
yalıtkan yüzey için görüntü yöntemiyle kapalı çözüm verir:
\begin{equation}
\varphi(r,z) \;=\; \frac{I}{4\pi\sqrt{\sigma_h\sigma_v}}
\left[\frac{1}{\sqrt{r^2+\Lambda^2 (z-z_s)^2}}
    + \frac{1}{\sqrt{r^2+\Lambda^2 (z+z_s)^2}}\right].
\label{eq:analytic}
\end{equation}
Bu ifade, izotropik limitte klasik $I/4\pi\sigma R$ çözümüne döner ve sayısal
çözücünün doğrulanmasında kullanılmıştır.

Ölçümün hangi derinliği gördüğü sorusu, direncin iletkenlik alanına göre
Fréchet türeviyle yanıtlanır. Karşılıklılık teoremi (Geselowitz), akım
çiftinin ürettiği $\varphi_A$ ile gerilim çiftinin birim akımla süründüğünde
üreteceği $\varphi_B$ alanları cinsinden
\begin{equation}
\frac{\delta R}{\delta \ln\sigma(\mathbf{x})}
\;=\; -\,\frac{\sigma(\mathbf{x})\,\nabla\varphi_A(\mathbf{x})\cdot
\nabla\varphi_B(\mathbf{x})}{I_A I_B}
\label{eq:sensitivity}
\end{equation}
verir. Bu çekirdeğin derinlik üzerine izdüşümü, elektrot diziliminin gerçek
araştırma derinliğini belirler ve Bölüm~\ref{sec:kernel}'de hesaplanmıştır.
Zaman serisi üretilirken her adımda tam potansiyel problemini çözmek yerine
birinci mertebe (Born) yaklaşımı kullanılır; bu yaklaşımın hatası aynı bölümde
tam çözümle karşılaştırılarak sayısallaştırılmıştır.

\subsection{İyi konumlanmışlık ve kararlılık}

Denklem~\eqref{eq:coupled} sisteminin parabolik olması, $D$ blokunun simetrik
kısmının pozitif tanımlı olmasını gerektirir. Tipik bir çalışma noktasında
($\theta = 0{,}15$, $T = 60\,^\circ$C) blok
\[
D \approx
\begin{pmatrix}
\dblkOO & \dblkOT\\[.2em]
\dblkTO & \dblkTT
\end{pmatrix}\ \mathrm{m^2 s^{-1}}
\]
değerini alır; simetrik kısmın en küçük özdeğeri $\coercivity$~m$^2$s$^{-1}$
olarak pozitiftir, yani sistem düzgün paraboliktir ve $L^2$ anlamında iyi
konumlanmıştır. Aynı blok, sistemin katılığını da açığa çıkarır: özdeğerler
arasındaki yaklaşık \stiffRatio\ kat fark, ısıl ve nem alt sistemlerinin çok
farklı zaman ölçeklerinde ilerlediğini gösterir.

Açık bir şema için von Neumann çözümlemesi, büyütme matrisi
$G(k) = \mathbf{I} - 4(\Delta t/\Delta z^2)\sin^2(k\Delta z/2)\,D$ üzerinden
$\rho(G)\le 1$ koşulunu $\Delta t \le \Delta z^2/2\lambda_{\max}(D)$ biçiminde
verir. Üretim ayarındaki $\Delta z = 2{,}5$~mm için bu sınır
$\dtExplicit$~s'dir. Çözücü bunun yerine yarı-örtük bir şema kullanır: difüzyon
terimleri örtük, katsayılar ve çapraz kaynaklar bir adım gecikmeli alınır. Bu,
doğrusal kararlılığı koşulsuz kılar ve üretim koşumlarının $\Delta t = 5$--$20$~s
ile ilerlemesine izin verir; bedeli zamanda birinci mertebe doğruluktur ve bu
bedelin büyüklüğü Bölüm~\ref{sec:conv}'de ölçülmüştür.

Bu tabloda gizli bir tuzak vardır ve çözümleme sırasında görülmüştür. Sağ üst
terim $L\rho_w D_{\theta v}/C$, yukarıdaki çalışma noktasında köşegen terimin
binde biridir; ancak $D_{\theta v}$ içindeki $\dd\psi/\dd\theta$ çarpanı, su
içeriği artık değere yaklaştığında van Genuchten bağıntısı gereği sınırsız
büyür. Kuru hücrelerde bu terim $10^{-4}$~m$^2$s$^{-1}$ mertebesine çıkar,
yani köşegen ısıl yayınımı iki mertebe aşar. Açık ayrıklaştırıldığında sonuç,
yüzey sıcaklığı sabitken bile adım başına onlarca kelvinlik salınımdır. Bu,
modelin değil şemanın kusurudur; Bölüm~\ref{sec:latent} bunun nasıl
kapatıldığını ve fiziksel bedelinin ne olduğunu verir.

% =============================================================================
\section{Ayrıklaştırma ve doğrulama}
\label{sec:conv}

\subsection{Şema}

Her iki denklem de sonlu hacim biçiminde, hücre merkezli ve yüz akılı olarak
ayrıklaştırılır. Yüz iletkenlikleri komşu hücre değerlerinin aritmetik
ortalamasıdır; iki boyutlu tensör operatöründe çapraz terimler, yüzü paylaşan
iki hücrenin merkezî farklarının ortalamasıyla kurulur, bu da düzgün ızgarada
ikinci mertebe verir. Yüzeyde sıcaklık için Dirichlet, nem için buharlaşma
akısı sınır koşulu uygulanır; tabanda sıcaklık sabit, nem serbest drenajlıdır.
Eksenel simetrik elektriksel problemde dış sınırlarda Dirichlet koşulu,
yarım hücre farkı yerine sınır yüzünde kurulan ikinci dereceden interpolant
üzerinden, $\partial_n\varphi|_b = (8\varphi_b - 9\varphi_N + \varphi_{N-1})/3h$
biçiminde kapatılır. Bu ayrıntı önemsiz değildir: yarım hücre kapanışı
kullanıldığında bütün şema, iç ayrıklaştırma ikinci mertebe olmasına rağmen
birinci mertebeye düşmektedir.

\subsection{Üretilmiş çözümlerle doğrulama}

İki uzaysal operatör de üretilmiş çözüm yöntemiyle (MMS) sınanmıştır. Isı
operatörü için, konuma bağlı ve çapraz bileşenli bir $\bK(x,z)$ ile
$T^\star = e^{-t/\tau}\cos(\pi x/L_x)\cos(\pi z/L_z)$ seçilmiş, karşılık gelen
kaynak analitik olarak türetilmiştir. Elektriksel operatör için
$\varphi^\star = \cos(ar)\cos(bz)$, $a=\pi/2r_{\max}$, $b=\pi/2z_{\max}$
seçimi yüzeyde sıfır akı ve dış sınırlarda sıfır potansiyel koşullarını tam
olarak sağlar ve tekillik içermez. Ölçülen yakınsama mertebeleri sırasıyla
$\pThermal$ ve $\pErt$'dir; ikisi de beklenen ikinci mertebeyi doğrular
(Şekil~\ref{fig:conv}a).

Aynı elektriksel çözücü, nokta kaynakla ve denklem~\eqref{eq:analytic}
analitik çözümüne karşı sınandığında mertebe $\pPoint$'e düşmektedir
(Şekil~\ref{fig:conv}b). Bu bir hata değil, birincil potansiyel biçiminde
$\delta$ kaynağının ayrıklaştırılmasının bilinen sonucudur; kaynak hücresindeki
tekillik küresel yakınsamayı birinci mertebeye sınırlar. Sonuç raporlanırken
ikinci mertebe varsayılmamış, doğrudan ölçülen mertebe kullanılmıştır.

\begin{figure}[htbp]\centering
\begin{tikzpicture}
\begin{groupplot}[group style={group size=2 by 1, horizontal sep=1.9cm},
  aegis, width=6.0cm, height=5.2cm, xmode=log, ymode=log]
\nextgroupplot[xlabel={Ağ adımı $h$}, ylabel={Bağıl $L_2$ hatası},
  legend pos=south east, title={(a) Üretilmiş çözümle doğrulama}]
\addplot[aegisteal, thick, mark=*, mark size=1.6] table[x=h,y=l2] {\dataroot/conv_mms_thermal.dat};
\addlegendentry{tensörlü ısı operatörü}
\addplot[aegisblue, thick, mark=square*, mark size=1.6] table[x=h,y=smooth] {\dataroot/conv_mms_ert.dat};
\addlegendentry{$\nabla\!\cdot\!(\sigma\nabla\phi)$}
\addplot[aegismuted, dashed, domain=0.004:0.06] {6e2*x^2};
\addlegendentry{$\mathcal{O}(h^2)$}
\nextgroupplot[xlabel={Ağ adımı $h$}, ylabel={Bağıl $L_2$ hatası},
  legend pos=south east, title={(b) Nokta kaynağın etkisi}]
\addplot[aegisember, thick, mark=triangle*, mark size=2] table[x=h,y=point] {\dataroot/conv_mms_ert.dat};
\addlegendentry{$\delta$ kaynaklı tam çözüm}
\addplot[aegismuted, dashed, domain=0.004:0.02] {0.32*x};
\addlegendentry{$\mathcal{O}(h)$}
\end{groupplot}
\end{tikzpicture}

\caption{Üretilmiş çözümlerle doğrulama. Solda iki uzaysal operatörün de
ikinci mertebeden yakınsadığı; sağda nokta kaynağın aynı çözücüde mertebeyi
birinciye düşürdüğü görülmektedir.}
\label{fig:conv}
\end{figure}

\subsection{İlgi büyüklüğünün ızgara ve zaman adımı yakınsaması}

Doğrulamanın operatör düzeyinde yapılması yeterli değildir; fizibilite kararı
öncü süreye dayandığından, yakınsama doğrudan bu büyüklük üzerinde de
ölçülmelidir. Öncü süre, eşik geçişlerinin örnekler arasında doğrusal ara
değerle bulunmasıyla sürekli bir büyüklüğe dönüştürülmüş, ardından ızgara
$\Delta z$ ve zaman adımı $\Delta t$ dört kademe inceltilmiştir
(Şekil~\ref{fig:gridconv}). Kaba ızgarada \leadCoarse~dk olan değer, en ince
ızgarada \leadFine~dk'ya oturur; gözlenen mertebe $\pGrid$ ve iki en ince
seviye arasındaki Izgara Yakınsama İndeksi \gciGrid~\%'dir. Zaman adımında
gözlenen mertebe $\pTime$'dir, yani beklendiği gibi birinci mertebeye yakındır.
Toplam sayısal belirsizlik, öncü süre üzerinde yaklaşık $\pm\numUnc$~dk
düzeyinde kalır; bu, aşağıdaki bölümlerde raporlanan fiziksel etkilerin çok
altındadır.

\begin{figure}[htbp]\centering
\begin{tikzpicture}
\begin{groupplot}[group style={group size=2 by 1, horizontal sep=1.8cm},
  aegis, width=5.6cm, height=4.9cm]
\nextgroupplot[xlabel={Hücre yüksekliği $\Delta z$ [mm]}, ylabel={Öncü süre [dk]},
  xmode=log, title={(a) Uzaysal ağ}]
\addplot[aegisteal, thick, mark=*, mark size=1.8] table[x=dz_mm,y=lead] {\dataroot/conv_grid.dat};
\nextgroupplot[xlabel={Zaman adımı $\Delta t$ [s]}, ylabel={Öncü süre [dk]},
  xmode=log, title={(b) Zaman adımı}]
\addplot[aegisember, thick, mark=square*, mark size=1.8] table[x=dt,y=lead] {\dataroot/conv_dt.dat};
\end{groupplot}
\end{tikzpicture}

\caption{Öncü sürenin ızgara ve zaman adımına göre yakınsaması. Değer, en ince
iki seviyede \gciGrid~\%'lik bir Izgara Yakınsama İndeksi bandına oturmaktadır.}
\label{fig:gridconv}
\end{figure}

\subsection{İzotermal gizli ısı teriminin ihmali}
\label{sec:latent}

Bir önceki bölümde işaret edilen kararsızlık, terimi sınırlayarak değil,
ihmalinin bedeli ölçülerek kapatılmıştır. Isı denklemindeki
$L\rho_w \Div(D_{\theta v}\nabla\theta)$ kaynağı, nem gradyanının sürüklediği
buharın taşıdığı gizli ısıdır ve sıcaklık gradyanının sürüklediği buharın
taşıdığı gizli ısıyla, yani $\lambda_{\mathrm{ef}}$ içindeki
$L\rho_w D_{Tv}$ terimiyle yarışır. Cephe civarında bu iki akının oranı
hesaplandığında ısıl olanın yaklaşık otuz kat büyük olduğu görülür. Kararlılığın
korunduğu bir zaman adımında ($\Delta t = 2$~s) model her iki biçimde de
koşturulmuş, öncü süre terim kapalıyken \latentOff, açıkken \latentOn~dakika
bulunmuştur; fark \latentDelta~\%'dir. Bu, aynı büyüklüğün ızgara ve zaman
adımı belirsizliğiyle karşılaştırılabilir düzeydedir. Terim bu nedenle ısı
denkleminden çıkarılmış, nem denkleminde ise örtük difüzyon katsayısı olarak
korunmuştur; orada işareti pozitif olduğu için kararlılık sorunu yoktur ve kuru
toprakta nem taşınımının baskın mekanizması zaten buhardır.

\subsection{Düzenlileştirme parametrelerine duyarsızlık}
\label{sec:reg}

Kaynama cephesinin kinetik hız katsayısı $k_0$ ve sıcaklık ölçeği $\Delta T$,
fiziksel sabitler değil sayısal düzenlileştirme parametreleridir. Dokuz
kombinasyon ($k_0 \in \{5\times10^{-4}, 10^{-3}, 5\times10^{-3}\}$~s$^{-1}$,
$\Delta T \in \{4,8,16\}$~K) için öncü süre \regLo\ ile \regHi~dk arasında,
yani \regSpread~\%'lik bir bant içinde kalmaktadır. Cephe bu nedenle sayısal
bir yapı değil, enerji dengesiyle belirlenen fiziksel bir yapı olarak kabul
edilebilir.

% =============================================================================
\section{Tespit fizibilitesi}

\subsection{Nominal senaryo}

Referans senaryo, ölü örtüde alevsiz pirolizle ilerleyen bir ısınmadır: yüzey
sıcaklığı bir buçuk saatte üç yüz dereceye yükselir, tutuşma yüzeyin bu eşiği
geçtiği an olarak tanımlanır. Prob altı santimetre derinliktedir.
Şekil~\ref{fig:nominal}, yüzey ve prob sıcaklıklarını, sıcaklık artış hızını
ve direncin logaritmik türevini birlikte gösterir. Tutuşma $t=\tIgnite$~saatte,
alarm $t=\tAlarm$~saatte gerçekleşir; öncü süre \leadNom~dakikadır.

İki nicelik bu grafikte dikkat çekicidir. Prob derinliğindeki sıcaklık artış
hızı en yüksek \maxdTdt~$^\circ$C/10~dk'ya ulaşır; oysa aynı derinlikte günlük
ısı dalgasının ürettiği doğal hız, sönümlü dalga çözümü
$T(z,t) = \bar T + A e^{-z/d}\sin(\omega t - z/d)$ ile
$d = \sqrt{2\alpha/\omega} = \dampDepth$~cm alınarak yalnızca
\diurnalRate~$^\circ$C/10~dk'dır. Belgelerdeki $\beta = 1{,}5$ eşiği bu doğal
arka planın \betaMarginNat\ katıdır, yani günlük döngü tek başına eşiği
tetikleyemez. İkinci nicelik direncin işaretidir ve bir sonraki alt bölümün
konusudur.

\begin{figure}[htbp]\centering
\begin{tikzpicture}
\begin{axis}[aegis, name=top, width=11.2cm, height=4.3cm, xmin=2.6, xmax=7,
  ylabel={Sıcaklık [$^\circ$C]}, legend pos=north west,
  xticklabels={}, ymin=0, ymax=460]
\addplot[aegisember, thick] table[x=t,y=Tsurf] {\dataroot/nominal_ts.dat};
\addlegendentry{yüzey $T(0,t)$}
\addplot[aegisteal, thick] table[x=t,y=Tprobe] {\dataroot/nominal_ts.dat};
\addlegendentry{prob $T(6\,\mathrm{cm},t)$}
\draw[aegisink, dashed, thin] (axis cs:4.5,0) -- (axis cs:4.5,460);
\node[font=\scriptsize, anchor=south east] at (axis cs:4.48,300) {tutuşma};
\end{axis}
\begin{axis}[aegis, at={(top.below south west)}, anchor=north west,
  width=11.2cm, height=4.3cm, xmin=2.6, xmax=7,
  xlabel={Zaman [saat]}, ylabel={$\mathrm{d}T/\mathrm{d}t$ [$^\circ$C/10\,dk]},
  legend pos=north east, ymin=-4, ymax=42]
\addplot[aegisteal, thick] table[x=t,y=dTdt] {\dataroot/nominal_ts.dat};
\addlegendentry{$\mathrm{d}T/\mathrm{d}t$}
\addplot[aegismuted, dashed, domain=2.6:7] {1.5}; \addlegendentry{$\beta=1{,}5$}
\end{axis}
\begin{axis}[aegis, grid=none, at={(top.below south west)}, anchor=north west,
  width=11.2cm, height=4.3cm, axis y line*=right, axis x line=none,
  xmin=2.6, xmax=7, ymin=-0.055, ymax=0.015, scaled y ticks=false,
  yticklabel style={aegisember, /pgf/number format/.cd, fixed, precision=3},
  ylabel={$\mathrm{d}(\ln R)/\mathrm{d}t$ [dk$^{-1}$]},
  ylabel style={aegisember}]
\addplot[aegisember, thick] table[x=t,y=dlnR] {\dataroot/nominal_ts.dat};
\end{axis}
\end{tikzpicture}

\caption{Nominal alev öncesi senaryo. Üstte yüzey ve prob sıcaklıkları; altta
sıcaklık artış hızı (sol eksen, $\beta$ eşiğiyle birlikte) ve direncin
logaritmik türevi (sağ eksen). Elektriksel imza, kuruma değil ısınma ve
yoğuşma kaynaklı olduğu için negatif yönde gelişmektedir.}
\label{fig:nominal}
\end{figure}

\subsection{Buharlaşma cephesi ve derinlik profilleri}

Şekil~\ref{fig:profile}, sıcaklık ve su içeriği profillerinin zaman içindeki
gelişimini verir. Sıcaklık profili yüz derece civarında belirgin bir düzlük
gösterir; bu, kaynama cephesinin gizli ısıyı soğurduğu bölgedir ve cephe
yavaşça derine ilerler. Su içeriği profili ise iki karşıt hareketi birlikte
sergiler: cephenin üstünde toprak hava kurusu hâle gelirken, cephenin hemen
altında su içeriği başlangıç değerinin üzerine çıkar. Bu yoğuşma bölgesi,
denklem~\eqref{eq:Dblock}'teki $D_{Tv}$ teriminin doğrudan sonucudur ve
ölçüm açısından belirleyicidir.

\begin{figure}[htbp]\centering
\begin{tikzpicture}
\begin{groupplot}[group style={group size=2 by 1, horizontal sep=1.7cm},
  aegis, width=5.2cm, height=6.0cm, y dir=reverse, ymin=0, ymax=32]
\nextgroupplot[xlabel={Sıcaklık [$^\circ$C]}, ylabel={Derinlik [cm]},
  legend pos=south east, xmin=0, xmax=340, title={(a) $T(z)$}]
\addplot table[x=T3p0,y=z] {\dataroot/nominal_profile.dat}; \addlegendentry{$t=3{,}0$ sa}
\addplot table[x=T4p0,y=z] {\dataroot/nominal_profile.dat}; \addlegendentry{$t=4{,}0$ sa}
\addplot table[x=T4p5,y=z] {\dataroot/nominal_profile.dat}; \addlegendentry{$t=4{,}5$ sa}
\addplot table[x=T6p0,y=z] {\dataroot/nominal_profile.dat}; \addlegendentry{$t=6{,}0$ sa}
\draw[aegisink, dashed, thin] (axis cs:100,0) -- (axis cs:100,32);
\nextgroupplot[xlabel={Hacimsel su içeriği $\theta$}, ylabel={Derinlik [cm]},
  legend pos=south east, xmin=0, xmax=0.30, title={(b) $\theta(z)$}]
\addplot table[x=th3p0,y=z] {\dataroot/nominal_profile.dat}; \addlegendentry{$t=3{,}0$ sa}
\addplot table[x=th4p0,y=z] {\dataroot/nominal_profile.dat}; \addlegendentry{$t=4{,}0$ sa}
\addplot table[x=th4p5,y=z] {\dataroot/nominal_profile.dat}; \addlegendentry{$t=4{,}5$ sa}
\addplot table[x=th6p0,y=z] {\dataroot/nominal_profile.dat}; \addlegendentry{$t=6{,}0$ sa}
\end{groupplot}
\end{tikzpicture}

\caption{Sıcaklık ve su içeriği profillerinin gelişimi. Yüz derecedeki düzlük
buharlaşma cephesidir; cephenin altında su içeriğinin başlangıç değerinin
üzerine çıkması, ısıl buhar sürüklenmesinin ürettiği yoğuşma bölgesidir.}
\label{fig:profile}
\end{figure}

\subsection{Elektriksel kuralın işareti: belgelerdeki tanımın düzeltilmesi}
\label{sec:sign}

Proje belgelerinde füzyon kuralının ikinci koşulu
$\dd(\ln R)/\dd t > \gamma$ biçimindedir, yani direncin \emph{artışını} arar.
Bağlaşık model, alev öncesi senaryoda bu koşulun geç kaldığını göstermektedir.
Isınma, gözenek suyunun iletkenliğini yüzde iki dereceye yakın bir katsayıyla
artırır; yoğuşma su içeriğini yükseltir; her iki etki de direnci düşürür.
Direnç ancak kuruma cephesi prob derinliğine ulaştığında hızla artar ve bu
cephe, ısıl cephenin belirgin biçimde gerisinden ilerler.
Nominal koşumda, altı santimetre derinlikte, tek yönlü kural tutuşmadan önce
hiç tetiklenmemekte; iki yönlü $|\dd(\ln R)/\dd t| > \gamma$ kuralı ise
\leadNom~dakika önce tetiklenmektedir. Aradaki fark, kuruma cephesinin prob
derinliğine varışını beklemekle, ısınmanın ve yoğuşmanın ürettiği erken imzayı
okumak arasındaki farktır.

Şekil~\ref{fig:depth} bu farkın prob derinliğine bağımlılığını verir. Tek yönlü
kural yalnızca prob \zOneMax~santimetreden sığ olduğunda alarm üretebilmekte,
en iyi durumda \zOneBest~santimetrede \leadOneBest~dakika bırakmaktadır; oysa
miselyumun yangın sırasında hayatta kaldığı bant on beş ile yirmi santimetre
arasındadır. Yani belgelerdeki kural, biyolojik olarak sürdürülemeyen bir
derinlik gerektirmektedir. İki yönlü kural ise altı santimetrede
\leadSix~dakika, on sekiz santimetrede \leadEighteen~dakika öncü süre üretir. Buradan çıkan tasarım önerisi, donanım raporundaki ayrıştırma
önerisiyle örtüşür ve onu güçlendirir: biyoelektrot miselyumun yaşadığı
derinlikte kalmalı, sıcaklık ve direnç probu aynı kablo üzerinde daha sığa
alınmalı, füzyon kuralının ikinci koşulu ise mutlak değer üzerinden
yazılmalıdır. Bu değişikliğin donanım maliyeti yoktur; yalnızca gömülü
yazılımdaki karşılaştırma işlecini ilgilendirir.

\begin{figure}[htbp]\centering
\begin{tikzpicture}
\begin{axis}[aegis, xlabel={Prob derinliği [cm]}, ylabel={Öncü süre [dk]},
  legend pos=north east, xmin=2, xmax=22, ymin=-15, ymax=105,
  width=10.4cm, height=5.4cm, restrict y to domain=-20:110]
\fill[aegisteal, opacity=.07] (axis cs:15,-15) rectangle (axis cs:20,105);
\node[font=\scriptsize, aegismuted, anchor=south, align=center]
  at (axis cs:17.5,-13) {miselyumun yaşadığı bant};
\addplot[aegisteal, thick, mark=*, mark size=1.5]
  table[x=z,y=lead_two] {\dataroot/depth_sweep.dat};
\addlegendentry{iki yönlü kural $|\mathrm{d}\ln R/\mathrm{d}t|>\gamma$}
\addplot[aegisember, thick, dashed, mark=square*, mark size=1.5]
  table[x=z,y=lead_one] {\dataroot/depth_sweep.dat};
\addlegendentry{belgedeki kural $\mathrm{d}\ln R/\mathrm{d}t>\gamma$}
\addplot[aegismuted, dotted, domain=2:22] {0};
\end{axis}
\end{tikzpicture}

\caption{Öncü sürenin prob derinliğine bağımlılığı, iki farklı $\gamma$ kuralı
için. Belgelerdeki tek yönlü kural yalnızca çok sığ problarda zamanında
çalışmakta, miselyumun yaşadığı bantta ise tutuşmadan sonra tetiklenmektedir.}
\label{fig:depth}
\end{figure}

\subsection{Ölçüm gerçekte hangi derinliği görüyor}
\label{sec:kernel}

Denklem~\eqref{eq:sensitivity} çekirdeği, düşey dizilimli dört elektrot için
(akım elektrotları üç ve dokuz, gerilim elektrotları beş ve yedi santimetrede)
hesaplanmıştır. Şekil~\ref{fig:kernel}, duyarlılık yoğunluğunu ve derinliğe
göre kümülatif dağılımını verir. Ölçümün medyan araştırma derinliği
\zMedian~santimetre, onuncu ve doksanıncı yüzdelikleri \zPten\ ile
\zPninety~santimetredir; yatayda duyarlılığın yüzde doksanı
\rPninety~santimetrelik bir yarıçap içinde kalır. Bu, elektrot çiftinin nokta
değil, belirli bir hacim üzerinde ağırlıklı ortalama ölçtüğünü ve bu hacmin
dizilim açıklığıyla ölçeklendiğini gösterir; dolayısıyla prob derinliğinin
değiştirilmesi, elektrot açıklığının da birlikte ölçeklenmesini gerektirir.

Zaman serisi üretiminde kullanılan birinci mertebe yaklaşım, aynı profiller
üzerinde tam eksenel simetrik çözümle karşılaştırılmıştır
(Şekil~\ref{fig:born}). Direncin logaritmik değişiminde en büyük bağıl sapma
\bornErr~\%'dir. Yaklaşımın Monte Carlo koşumlarında kullanılması, her adımda
seyrek bir doğrusal sistem çözmeye göre üç mertebe hızlıdır ve bu sapma
düzeyinde karar sınırını değiştirmemektedir.

\begin{figure}[htbp]\centering
\begin{tikzpicture}
\begin{groupplot}[group style={group size=2 by 1, horizontal sep=1.8cm},
  aegis, width=5.4cm, height=5.2cm, xmin=0, xmax=30]
\nextgroupplot[xlabel={Derinlik [cm]}, ylabel={Duyarlılık yoğunluğu},
  title={(a) Geselowitz çekirdeği}]
\addplot[aegisteal, thick, fill=aegisteal, fill opacity=.18]
  table[x=z,y=w] {\dataroot/ert_kernel.dat} \closedcycle;
\nextgroupplot[xlabel={Derinlik [cm]}, ylabel={Kümülatif duyarlılık},
  ymin=0, ymax=1.02, title={(b) Araştırma derinliği}]
\addplot[aegisblue, thick] table[x=z,y=cum] {\dataroot/ert_kernel.dat};
\addplot[aegismuted, dashed, domain=0:30] {0.5};
\addplot[aegismuted, dotted, domain=0:30] {0.9};
\end{groupplot}
\end{tikzpicture}

\caption{Dört elektrotlu düşey dizilimin Geselowitz duyarlılık çekirdeği ve
kümülatif araştırma derinliği. Ölçüm, medyanda \zMedian~cm derinliği
görmektedir.}
\label{fig:kernel}
\end{figure}

\begin{figure}[htbp]\centering
\begin{tikzpicture}
\begin{axis}[aegis, xlabel={Zaman [saat]}, ylabel={$\ln\!\big(R(t)/R(t_0)\big)$},
  legend pos=south west, width=9.6cm, height=5.0cm]
\addplot[aegisteal, thick, mark=*, mark size=1.3] table[x=t,y=lnfull] {\dataroot/born_check.dat};
\addlegendentry{tam eksenel simetrik çözüm}
\addplot[aegisember, dashed, thick, mark=o, mark size=1.6] table[x=t,y=lnborn] {\dataroot/born_check.dat};
\addlegendentry{birinci mertebe (Born) yaklaşımı}
\end{axis}
\end{tikzpicture}

\caption{Birinci mertebe (Born) yaklaşımının tam eksenel simetrik çözümle
karşılaştırılması. En büyük bağıl sapma \bornErr~\%'dir.}
\label{fig:born}
\end{figure}

\subsection{Belirsizlik altında öncü süre}

Tek bir nominal koşum fizibilite kanıtı değildir. Sekiz parametreli bir
tasarım uzayında (başlangıç nem içeriği, ortam sıcaklığı, piroliz ramp süresi,
yüzey plato sıcaklığı, doygun ısıl iletkenlik, Archie üssü, prob derinliği ve
gözenek suyu iletkenliği) \mcPos\ örneklik bir Sobol dizisi çekilmiş ve tümü
yığın hâlinde çözülmüştür. Olayların \detRate~\%'sinde alarm üretilmiştir;
öncü sürenin medyanı \leadMed~dakika, beşinci ve doksan beşinci yüzdelikleri
\leadPfive\ ile \leadPnf~dakikadır. Onbeş dakikalık asgari müdahale eşiğinin
üzerinde kalma olasılığı \pLeadFifteen~\%, otuz dakikanın üzerinde kalma
olasılığı \pLeadThirty~\%'dir (Şekil~\ref{fig:leadcdf}).

\begin{figure}[htbp]\centering
\begin{tikzpicture}
\begin{axis}[aegis, xlabel={Öncü süre [dk]}, ylabel={Birikimli olasılık},
  xmin=0, xmax=150, ymin=0, ymax=1.02, width=9.6cm, height=5.2cm,
  legend pos=south east]
\addplot[aegisteal, very thick] table[x=lead,y=cdf] {\dataroot/lead_cdf.dat};
\addlegendentry{Monte Carlo ampirik dağılım}
\addplot[aegismuted, dashed, domain=0:150] {0.05};
\addplot[aegismuted, dashed, domain=0:150] {0.95};
\draw[aegisember, thin, dash dot] (axis cs:15,0) -- (axis cs:15,1);
\node[font=\scriptsize, aegisember, anchor=south west] at (axis cs:16,0.04)
  {müdahale eşiği};
\end{axis}
\end{tikzpicture}

\caption{Öncü sürenin belirsizlik altındaki ampirik birikimli dağılımı.
Kesikli çizgiler beşinci ve doksan beşinci yüzdelikleri işaretlemektedir.}
\label{fig:leadcdf}
\end{figure}

\subsection{Ayırt edicilik ve eşik tasarım penceresi}

Fizibilitenin ikinci sorusu, bu imzanın yangın olmadan da gerçekleşen
olaylardan ayrılabilirliğidir. Aynı parametre uzayında kuraklık ve normal gün
senaryoları için toplam \mcNeg\ örneklik olumsuz kümeler üretilmiş, karar
istatistiği olarak normalleştirilmiş füzyon büyüklüğü
$\Lambda = \max_t \min\!\left(\dot T/\beta_0,\;|\dd\ln R/\dd t|/\gamma_0\right)$
tanımlanmıştır.

Sonuç, beklenenden keskindir. Her iki fiziksel kanal da bu senaryo aileleri
altında olumlu ve olumsuz kümeleri örtüşmesiz biçimde ayırmakta, eğri altı
alanları \aucFusion, \aucBeta\ ve \aucGamma\ çıkmaktadır; belgelerdeki eşik
çiftine karşılık gelen $\tau = 1$ çalışma noktasında doğru tespit oranı
\tprOne~\%, yanlış alarm oranı \fprOne~\%'dir. Alıcı işletim eğrisi bu durumda
köşeye yapıştığı için bilgi taşımaz; ayrışmanın gerçek ölçüsü, iki kümenin
dağılımları arasındaki boşluktur (Şekil~\ref{fig:separation}).

Bu boşluk, eşiklerin tasarım penceresini doğrudan verir. Isıl kanalda olumsuz
kümenin üst kuyruğu ile alev öncesi kümenin alt kuyruğu arasında
$\beta \in [\betaLo, \betaHi]$, elektriksel kanalda
$\gamma \in [\gammaLo, \gammaHi]$ bandı kalmaktadır; belgelerdeki
$\beta = 1{,}5$ ve $\gamma = 0{,}002$ değerleri her iki pencerenin de içindedir.

Kuraklık senaryosunun ayrıntısı, mantıksal çarpımın neden yine de doğru tasarım
olduğunu gösterir. Kuraklıkta elektriksel kanal en büyük \droughtdlnR~dk$^{-1}$,
ısıl kanal \droughtdT~$^\circ$C/10~dk değerine ulaşır; alev öncesi senaryoda
ısıl kanal \maxdTdt~$^\circ$C/10~dk'ya, yani kuraklıktakinin yaklaşık \sepBeta\
katına çıkar. Modellenen senaryo ailesinde tek bir kanal da yeterlidir; ancak
bu, kanalın kendisinin arızalanmadığı, sürüklenmediği ve yerel bir bozucudan
etkilenmediği varsayımına dayanır. Çarpımın değeri, ölçülebilir bir eğri altı
alanı kazancında değil, tek kanallı arıza kiplerine karşı sağladığı
yedeklilikte aranmalıdır; bu da modelin içinden değil saha verisinden
doğrulanabilecek bir iddiadır.

\begin{figure}[htbp]\centering
\begin{tikzpicture}
\begin{groupplot}[group style={group size=2 by 1, horizontal sep=1.9cm},
  aegis, width=5.6cm, height=5.2cm, xmode=log, ymin=0, ymax=1.02,
  ylabel={Birikimli olasılık}]
\nextgroupplot[xlabel={$\max_t \mathrm{d}T/\mathrm{d}t$ [$^\circ$C/10\,dk]},
  legend pos=north west, title={(a) Isıl kanal}, xmin=0.05, xmax=200]
\addplot[aegisblue, thick] table[x=beta_neg, y expr=\coordindex/1023] {\dataroot/threshold_box.dat};
\addlegendentry{olumsuz küme}
\addplot[aegisteal, thick] table[x=beta_pos, y expr=\coordindex/1023] {\dataroot/threshold_box.dat};
\addlegendentry{alev öncesi}
\draw[aegisember, thick, dashed] (axis cs:1.5,0) -- (axis cs:1.5,1.02);
\node[font=\scriptsize, aegisember, anchor=south west, rotate=90]
  at (axis cs:1.6,0.06) {$\beta = 1{,}5$};
\nextgroupplot[xlabel={$\max_t |\mathrm{d}\ln R/\mathrm{d}t|$ [dk$^{-1}$]},
  legend pos=north west, title={(b) Elektriksel kanal}, xmin=1e-4, xmax=1]
\addplot[aegisblue, thick] table[x=gamma_neg, y expr=\coordindex/1023] {\dataroot/threshold_box.dat};
\addlegendentry{olumsuz küme}
\addplot[aegisteal, thick] table[x=gamma_pos, y expr=\coordindex/1023] {\dataroot/threshold_box.dat};
\addlegendentry{alev öncesi}
\draw[aegisember, thick, dashed] (axis cs:0.002,0) -- (axis cs:0.002,1.02);
\node[font=\scriptsize, aegisember, anchor=south west, rotate=90]
  at (axis cs:0.0022,0.06) {$\gamma = 0{,}002$};
\end{groupplot}
\end{tikzpicture}

\caption{İki kanalın olumlu ve olumsuz kümelerdeki birikimli dağılımları.
Kesikli çizgiler belgelerdeki eşikleri işaretler; eşik ile her iki dağılımın
kuyruğu arasındaki boşluk, tasarım marjını doğrudan görünür kılar.}
\label{fig:separation}
\end{figure}

\subsection{Küresel duyarlılık}

Hangi belirsizliğin sonucu belirlediği, Saltelli tasarımıyla \sobolRuns\
koşumda hesaplanan Sobol indisleriyle yanıtlanır
(Şekil~\ref{fig:sobol}). Toplam etki bakımından baskın parametre
\sobolFirst\ ($S_T = \sobolFirstST$), ardından \sobolSecond\
($S_T = \sobolSecondST$) ve \sobolThird\ ($S_T = \sobolThirdST$)
gelmektedir. Birinci mertebe indislerin toplamı \sobolSum{}'dir; birden
belirgin biçimde küçük olması, parametreler arasında etkileşimin önemli
olduğunu, dolayısıyla tek tek duyarlılık taramalarının yanıltıcı olacağını
gösterir.

Çözümlemenin kendi doğruluğu için tasarıma bilinçli olarak etkisiz bir
parametre eklenmiştir. Gözenek suyu iletkenliği $\sigma_w^{25}$,
denklem~\eqref{eq:archie}'de çarpımsal bir sabit olduğundan direncin
logaritmik türevinden düşer ve öncü süreyi etkileyemez. Hesaplanan toplam
etki indisi \sobolNull{}'dur; tahmin edici, etkisiz parametreyi doğru biçimde
sıfırda vermektedir.

Aynı grafikte elektriksel kanalın diğer parametresi olan Archie üssünün de
sıfır çıkması, sayısal bir tesadüf değil yapısal bir sonuçtur ve tasarım
açısından önemlidir. İki yönlü kural altında elektriksel koşul, ısıl koşuldan
\emph{daha erken} sağlanmaktadır; dolayısıyla alarm anını belirleyen, iki
koşuldan geç olanı, yani sıcaklık kanalıdır. Elektriksel kanalın işlevi öncü
süreyi uzatmak değil, kuraklığı yangından ayırmaktır. Bu iş bölümü, bir önceki
bölümdeki ayırt edicilik sonuçlarıyla birlikte okunmalıdır: sıcaklık kanalı
zamanlamayı, direnç kanalı özgüllüğü taşır.

\begin{figure}[htbp]\centering
\begin{tikzpicture}
\begin{axis}[aegis, ybar, bar width=7pt, width=10.4cm, height=5.0cm,
  ylabel={Sobol indisi},
  xtick=data, xticklabels={$\theta_0$,$T_\infty$,$t_r$,$T_{\max}$,$\lambda_{\mathrm{sat}}$,$n_A$,$z_p$,$\sigma_w$},
  x tick label style={font=\small}, ymin=0, legend pos=north east,
  enlarge x limits=0.07]
\addplot[fill=aegisteal, draw=aegisdeep] table[x=idx,y=S1] {\dataroot/sobol.dat};
\addlegendentry{birinci mertebe $S_i$}
\addplot[fill=aegisember!70, draw=aegisember] table[x=idx,y=ST] {\dataroot/sobol.dat};
\addlegendentry{toplam etki $S_{T_i}$}
\end{axis}
\end{tikzpicture}

\caption{Öncü süre üzerindeki birinci mertebe ve toplam etki Sobol indisleri.
Sağdaki $\sigma_w$ sütunu, kuramsal olarak etkisiz olması gereken kontrol
parametresidir.}
\label{fig:sobol}
\end{figure}

% =============================================================================
\section{Mekânsal fizibilite}
\label{sec:lateral}

\subsection{Yanal ısı yayılımı ve eğimin etkisi}

Buraya kadarki çözümleme düşey bir sütun üzerindedir ve bir düğümün tam
altında başlayan bir olayı varsayar. Gerçek soru, düğümden belirli bir yatay
uzaklıkta başlayan bir sıcak noktanın ne zaman görüleceğidir. Tam tensörlü iki
boyutlu çözücü, yüzeyde \latSpot~metre yarıçaplı bir sıcak nokta ile
koşturulmuş ve altı santimetre derinlikte $\beta$ eşiğinin aşıldığı an, yatay
uzaklığın fonksiyonu olarak çıkarılmıştır (Şekil~\ref{fig:lateral}).

Sonuç, ısıl kanalın mekânsal erimi konusunda serinkanlı olmayı gerektirir.
Yatay yatakta tespit yarıçapı \latRzero~metredir; ısı iletiminin difüzyon
uzunluğu $\sqrt{4\alpha t}$ ile ölçeklendiği düşünülürse bu beklenen bir
büyüklüktür. Yatak eğimli olduğunda ayak izi simetrisini yitirir: bir yönde
\latRdown, karşı yönde \latRup~metre — küçük ama tutarlı bir fark. Bu asimetri,
Şekil~\ref{fig:tensor}'deki akı sapmasının mekânsal karşılığıdır ve eğimli
parsellerde düğüm yerleşiminin yamaç yönüne göre kaydırılabileceğini
düşündürür. Pratik
sonucu şudur: ısıl kanal tek başına metre mertebesinde bir erime sahiptir ve
hektarlar ölçeğinde kapsama sağlayamaz. Sistemin mekânsal kapsama iddiası bu
nedenle ısıl kanala değil, biyolojik kanala dayanmak zorundadır.

\begin{figure}[htbp]\centering
\begin{tikzpicture}
\begin{axis}[aegis, xlabel={Sıcak nokta merkezine yatay uzaklık [m]},
  ylabel={$\beta$ eşiğinin aşıldığı an [dk]}, xmin=-2.5, xmax=2.5,
  width=9.6cm, height=5.0cm, legend pos=north east,
  restrict y to domain=150:360]
\addplot[aegisteal, thick] table[x=x,y=t_dip0] {\dataroot/lateral.dat};
\addlegendentry{yatay yatak ($\psi=0^\circ$)}
\addplot[aegisember, thick, dashed] table[x=x,y=t_dip25] {\dataroot/lateral.dat};
\addlegendentry{eğimli yatak ($\psi=25^\circ$)}
\end{axis}
\end{tikzpicture}

\caption{İki boyutlu tensörlü çözümde yanal tespit ayak izi. Eğimli yatakta
ayak izi yamaç yönünde kayarak asimetrikleşmektedir.}
\label{fig:lateral}
\end{figure}

\subsection{Miselyum ağının kablo denklemi}

Miselyum ağı, elektriksel açıdan dağıtılmış bir iletken ortamdır ve zar
üzerinden sızıntı ile eksen boyunca iletimin dengelendiği pasif kablo
denklemiyle temsil edilebilir. Ağ toprakta yatay bir tabaka gibi yayıldığından
iki boyutlu kararlı hâl biçimi uygundur:
\begin{equation}
\nabla^2 V - \frac{V}{\lambda^2} = 0,
\qquad \lambda = \sqrt{\frac{r_m}{r_i}},
\end{equation}
burada $\lambda$ elektrotonik uzunluk sabitidir. Yarıçapı $a$ olan dairesel
bir stres kaynağı için çözüm, sıfırıncı mertebe değiştirilmiş Bessel
fonksiyonuyla verilir:
\begin{equation}
V(r) = V_0\,\frac{K_0(r/\lambda)}{K_0(a/\lambda)} .
\label{eq:cable}
\end{equation}
Deney raporlarındaki genlikler kaynak için $V_0 = \cableVzero$~mV mertebesini,
ölçüm zincirinin gürültü tabanı ve ortam değişkenliği ise
$V_{\mathrm{det}} = \cableVdet$~mV'lik bir tespit eşiğini makul kılar.
Şekil~\ref{fig:cable}, farklı $\lambda$ değerleri için sönümü ve karşılık gelen
tespit yarıçapını verir. Elde edilen ilişki, katsayısı $\lambda$ büyüdükçe
hafifçe düşen bir doğrusallıktır: $\lambda = 1$~m için \rdetOne~m,
$\lambda = 10$~m için \rdetTen~m, $\lambda = 30$~m için \rdetThirty~m.

Bu, fizibilitenin en kritik bilinmeyenidir. Ağın $\lambda$ değeri hiçbir proje
belgesinde ölçülmemiştir ve literatürde de doğrudan bir değeri yoktur.
Aşağıdaki alt bölüm, bu bilinmeyeni ekonomik bir eşiğe çevirerek ölçülebilir
bir hedefe dönüştürür.

\begin{figure}[htbp]\centering
\begin{tikzpicture}
\begin{groupplot}[group style={group size=2 by 1, horizontal sep=1.9cm},
  aegis, width=5.4cm, height=5.2cm]
\nextgroupplot[xlabel={Ağ boyunca uzaklık $r$ [m]}, ylabel={Genlik [mV]},
  xmode=log, ymode=log, xmin=0.05, xmax=60, ymin=1e-3, ymax=60,
  legend pos=south west, title={(a) $V(r)=V_0K_0(r/\lambda)/K_0(a/\lambda)$}]
\addplot table[x=r,y=V0p5] {\dataroot/cable.dat}; \addlegendentry{$\lambda=0{,}5$ m}
\addplot table[x=r,y=V1]   {\dataroot/cable.dat}; \addlegendentry{$\lambda=1$ m}
\addplot table[x=r,y=V10]  {\dataroot/cable.dat}; \addlegendentry{$\lambda=10$ m}
\addplot table[x=r,y=V30]  {\dataroot/cable.dat}; \addlegendentry{$\lambda=30$ m}
\addplot[aegisink, dashed, domain=0.05:60] {0.5};
\nextgroupplot[xlabel={Elektrotonik uzunluk $\lambda$ [m]},
  ylabel={Tespit yarıçapı $r_{\mathrm{det}}$ [m]},
  xmode=log, ymode=log, title={(b) Kapsama yarıçapı}]
\addplot[aegisteal, thick, mark=*, mark size=1.8] table[x=lam,y=rdet] {\dataroot/cable_rdet.dat};
\end{groupplot}
\end{tikzpicture}

\caption{İki boyutlu kablo denkleminin çözümü ve tespit yarıçapının
elektrotonik uzunluk sabitine bağımlılığı.}
\label{fig:cable}
\end{figure}

\subsection{Düğüm aralığı, kapsama ve maliyet}

Düğümler altıgen örgüde $S$ aralığıyla yerleştirilirse hektardaki düğüm sayısı
$10^4/(0{,}866 S^2)$, herhangi bir noktanın en yakın düğüme uzaklığının en
büyük değeri $S/\sqrt{3}$ olur. Bir tutuşma, düğümün tespit yarıçapı içinde
başlarsa alev öncesi imza görülür; dışında başlarsa yangının yayılarak bu
yarıçapa girmesi beklenir. İkinci durumda tespit anındaki yanmış alan yalnızca
geometriye bağlıdır:
\begin{equation}
A_{\mathrm{tespit}} = \pi\left(\frac{S}{\sqrt3} - r_{\mathrm{det}}\right)^{\!2},
\qquad
t_{\mathrm{tespit}} = \frac{S/\sqrt3 - r_{\mathrm{det}}}{v},
\end{equation}
burada $v$ yüzey yangınının yayılma hızıdır. Alev öncesi tespit olasılığı ise
$\pi r_{\mathrm{det}}^2/(0{,}866 S^2)$ oranıdır.

Şekil~\ref{fig:coverage}, bu üç büyüklüğü düğüm aralığına göre verir. Yüz
metrelik aralık, hektar başına \covNodes\ düğüm ve \covCapex~TL kurulum
maliyeti demektir; bu yoğunlukta en kötü durumda tespit, yangın
\covArea~hektara ulaştığında, tutuşmadan yaklaşık \covTime~dakika sonra
gerçekleşir. Aynı yoğunlukta alev öncesi tespit olasılığı $\lambda$'ya
kuvvetle bağlıdır; $\lambda$ bir metre mertebesindeyse yüzde birkaç, on metre
mertebesindeyse yüzde yirmilere çıkar. Bu nedenle sistemin ayırt edici
değer önerisi, düğüm başına maliyeti düşürmekten çok $\lambda$'yı ölçmek ve
mümkünse elektrot yüzeyi ile inokülasyon protokolü üzerinden büyütmekle
ilgilidir.

Dürüst bir fizibilite ifadesi şudur: bu maliyet düzeyi, bir orman alanının
tamamını battaniye gibi örtmek için uygun değildir; enerji ve maden sahaları,
yangın emniyet şeritleri ve yerleşim--orman arayüzü gibi yüksek değerli ve
sınırlı geometriler için uygundur. İş planındaki aşamalı pazar sıralaması
zaten bu geometrileri hedeflemektedir ve bu çözümleme onu sayısal olarak
desteklemektedir.

\begin{figure}[htbp]\centering
\begin{tikzpicture}
\begin{axis}[aegis, xlabel={Düğüm aralığı $S$ [m]},
  ylabel={Alev öncesi tespit olasılığı}, xmin=20, xmax=200, ymin=0, ymax=1,
  width=9.4cm, height=5.2cm, legend pos=north east]
\addplot table[x=S,y=p1]  {\dataroot/coverage.dat}; \addlegendentry{$\lambda=1$ m}
\addplot table[x=S,y=p3]  {\dataroot/coverage.dat}; \addlegendentry{$\lambda=3$ m}
\addplot table[x=S,y=p10] {\dataroot/coverage.dat}; \addlegendentry{$\lambda=10$ m}
\addplot table[x=S,y=p30] {\dataroot/coverage.dat}; \addlegendentry{$\lambda=30$ m}
\end{axis}
\begin{axis}[aegis, grid=none, width=9.4cm, height=5.2cm,
  axis y line*=right, axis x line=none, xmin=20, xmax=200,
  ymode=log, ymin=3e3, ymax=2e6, ylabel={Kurulum maliyeti [TL/ha]},
  ylabel style={aegisember}, y tick label style={aegisember}]
\addplot[aegisember, very thick] table[x=S,y=capex] {\dataroot/coverage.dat};
\end{axis}
\end{tikzpicture}

\caption{Alev öncesi tespit olasılığı (sol eksen, farklı elektrotonik uzunluk
sabitleri için) ve hektar başına kurulum maliyeti (sağ eksen), düğüm aralığına
göre.}
\label{fig:coverage}
\end{figure}

% =============================================================================
\section{Coğrafi fizibilite: İzmir Orman Bölge Müdürlüğü kapsama alanı}

\subsection{Veri, izdüşüm ve türetilen alanlar}

Modelin gerçek arazide ne ürettiğini görmek için çözümleme, 26,5--28,2 doğu
boylamları ile 38,0--39,0 kuzey enlemleri arasındaki kutuya taşınmıştır. SRTM
bir yay-saniyelik sayısal yükseklik modeli mozaiklenmiş, UTM 35K izdüşümüne
\geoRes~metre çözünürlükte yeniden örneklenmiş, kara alanı \geoLandArea~km$^2$,
en yüksek nokta \geoElevMax~metre olarak bulunmuştur. Eğim ve bakı Horn'un
üç-üçlük operatörüyle, akış birikimi öncelikli-taşkın algoritmasıyla çukurları
doldurulmuş yüzey üzerinde D8 yöntemiyle, topografik ıslaklık indisi
$\mathrm{TWI} = \ln(a/\tan\beta)$ bağıntısıyla hesaplanmıştır. Eğimli yüzeye
düşen potansiyel doğrudan güneş ışınımı, gün boyunca on dakikalık adımlarla
gökyüzü geometrisi integre edilerek elde edilmiştir. İl sınırları, kıyı çizgisi
ve yerleşim noktaları Natural Earth veri kümesinden okunmuş, aynı izdüşüme
alınmış ve yerleşim--orman arayüzü mesafesi bir $k$-boyutlu ağaç sorgusuyla
çıkarılmıştır. Kapsama alanının ortalama eğimi \geoSlopeMean\ derecedir.

\subsection{Tutuşma duyarlılığı ve öncü sürenin mekânsal dağılımı}

Çok ölçütlü duyarlılık indisi, Rothermel'in eğim çarpanıyla orantılı bir
yayılma terimi, normalize potansiyel ışınım, ıslaklık indisinin tümleyeni ve
yerleşime yakınlık teriminin ağırlıklı toplamıdır. İndisin altmış ikinci
yüzdeliğin üzerinde kaldığı alan \geoHighArea~km$^2$'dir
(Şekil~\ref{fig:georisk}a).

Asıl bağ burada kurulur. Arazi eğimi $\psi$, denklem~\eqref{eq:Ktensor}
tensöründe yatak normalini döndüren açının ta kendisidir; düşey iletim bileşeni
$K_{zz} = k_\parallel\sin^2\psi + k_\perp\cos^2\psi$ olarak eğimle
büyür. Islaklık indisi başlangıç su içeriğini, yükseklik ve ışınım ortam
sıcaklığını belirler. Bu üç alan, arazi hücrelerinden örneklenen \geoCells\
noktada bağlaşık sütun modeline girdi olarak verilmiş ve model yığın hâlinde
çözülmüştür. Elde edilen öncü süre alanının medyanı \geoLeadMed~dakika,
beşinci ve doksan beşinci yüzdelikleri \geoLeadPfive\ ile \geoLeadPnf~dakikadır
(Şekil~\ref{fig:georisk}b). Dağılımın genişliği, tek bir nominal koşumun
saha için yeterli olmadığını; ancak alt kuyruğun bile müdahale eşiğinin
üzerinde kaldığını göstermektedir.

\begin{figure}[htbp]\centering
\begin{tikzpicture}
\begin{axis}[aegismap, width=10.4cm, height=6.6cm,
  colorbar, colormap name=aegisterrain, point meta min=-120, point meta max=1500,
  colorbar style={ylabel={Yükseklik [m]}},
  unbounded coords=jump]
\addplot[matrix plot*, point meta=explicit, mesh/cols=139]
  table[x=x,y=y,meta=z] {\dataroot/geo_dem.dat};
\addplot[aegismuted, line width=.25pt, opacity=.55] table[x=x,y=y] {\dataroot/geo_graticule.dat};
\addplot[aegisink, line width=.5pt] table[x=x,y=y] {\dataroot/geo_coast.dat};
\addplot[aegismuted, line width=.35pt, dashed] table[x=x,y=y] {\dataroot/geo_prov_lines.dat};
\addplot[only marks, mark=*, mark size=2.0, aegisember,
  mark options={draw=white, line width=.4pt}]
  table[x=x,y=y] {\dataroot/geo_sites.dat};
\end{axis}
\end{tikzpicture}

\caption{Kapsama alanının sayısal yükseklik modeli; üzerine kıyı çizgisi, il
sınırları ve yarım derecelik enlem--boylam ağı bindirilmiştir. Ağ çizgilerinin
düz değil hafifçe eğri olması, UTM izdüşümünün kendi geometrisidir. Turuncu
noktalar, eğim, yükseklik ve yerleşim mesafesi kısıtları altında en yüksek
duyarlılık indisine sahip pilot saha adaylarıdır.}
\label{fig:geomap}
\end{figure}

\begin{figure}[htbp]\centering
\begin{tikzpicture}
\begin{groupplot}[group style={group size=2 by 1, horizontal sep=2.8cm},
  aegismap, width=5.0cm, height=4.4cm, unbounded coords=jump]
\nextgroupplot[colorbar, colormap/hot, point meta min=-0.06, point meta max=0.85,
  colorbar style={ylabel={indis}}, title={(a) Tutuşma duyarlılığı}]
\addplot[matrix plot*, point meta=explicit, mesh/cols=139]
  table[x=x,y=y,meta=risk] {\dataroot/geo_dem.dat};
\addplot[aegisink, line width=.4pt] table[x=x,y=y] {\dataroot/geo_coast.dat};
\nextgroupplot[colorbar, colormap/viridis,
  colorbar style={ylabel={öncü süre [dk]}}, title={(b) Modellenen öncü süre}]
\addplot[scatter, only marks, mark=square*, mark size=.55,
  scatter src=explicit, draw opacity=0]
  table[x=x,y=y,meta=lead] {\dataroot/geo_lead_pts.dat};
\addplot[aegisink, line width=.4pt] table[x=x,y=y] {\dataroot/geo_coast.dat};
\end{groupplot}
\end{tikzpicture}

\caption{Solda çok ölçütlü tutuşma duyarlılık indisi; sağda aynı arazi
üzerinde bağlaşık modelin ürettiği öncü süre alanı.}
\label{fig:georisk}
\end{figure}

\subsection{İl bazında dağılım ve pilot saha adayları}

Aynı ızgara, il sınırlarıyla maskelenerek \geoNprov\ il için özetlenmiştir
(Şekil~\ref{fig:geoprov}). Ortalama duyarlılık bakımından ilk sırada
\geoTopProv\ (\geoTopRisk) yer almaktadır. Pilot saha seçimi, indis
sıralamasına üç kısıt eklenerek yapılmıştır: eğim otuz iki dereceden küçük,
yükseklik bin üç yüz metreden az ve yerleşime uzaklık sekiz yüz metre ile on
iki kilometre arasında. Son kısıt, hem erişim lojistiği hem de arayüz
yangınlarının fiilen gerçekleştiği bant içindir. Birbirinden en az üç
kilometre uzak \geoSiteN\ aday elde edilmiştir; en yüksek puanlı adayın indisi
\geoSiteRisk, yüksekliği \geoSiteElev~metre, eğimi \geoSiteSlope\ derecedir
(Şekil~\ref{fig:geomap}).

\begin{figure}[htbp]\centering
\begin{tikzpicture}
\begin{axis}[aegis, ybar, bar width=9pt, width=9.4cm, height=4.6cm,
  ylabel={Ortalama duyarlılık indisi}, xtick=data,
  xticklabel style={font=\scriptsize, rotate=30, anchor=north east},
  ymin=0, enlarge x limits=0.12, legend pos=north east]
\addplot[fill=aegisember!55, draw=aegisember] table[x=idx,y=risk_mean] {\dataroot/geo_prov.dat};
\addlegendentry{ortalama}
\addplot[fill=aegisteal!45, draw=aegisdeep] table[x=idx,y=risk_p90] {\dataroot/geo_prov.dat};
\addlegendentry{90. yüzdelik}
\end{axis}
\end{tikzpicture}

\caption{Kapsama alanındaki iller için ortalama ve doksanıncı yüzdelik
duyarlılık indisi.}
\label{fig:geoprov}
\end{figure}

\subsection{Kesit: arazi ve yeraltı birlikte}

Şekil~\ref{fig:geosection}, 38,42 kuzey enlemi boyunca uzanan
\geoTransLen~kilometrelik bir batı--doğu kesitini verir. Üstteki panel arazi
profilini, alttaki panel aynı kesit boyunca bağlaşık modelin beşinci saatte
ürettiği yeraltı sıcaklık alanını gösterir. Kesit boyunca öncü sürenin medyanı
\geoTransLead\ dakikadır. Yüz derecelik buharlaşma cephesinin derinliği kesit
boyunca değişmekte, kuru ve dik yamaçlarda daha derine inmektedir; bu, aynı
donanımın farklı arazi birimlerinde farklı öncü süreler üreteceğinin ve pilot
saha seçiminin yalnızca risk değil, model duyarlılığı açısından da anlamlı
olduğunun görsel kanıtıdır.

\begin{figure}[htbp]\centering
\begin{tikzpicture}
\begin{axis}[aegis, name=prof, width=10.6cm, height=3.0cm, xmin=0, xmax=126,
  ylabel={Yükseklik [m]}, xticklabels={}, ymin=0]
\addplot[aegisteal, thick, fill=aegisteal!18] table[x=s,y=z] {\dataroot/geo_transect.dat}
  \closedcycle;
\end{axis}
\begin{axis}[aegis, at={(prof.below south west)}, anchor=north west,
  width=10.6cm, height=4.0cm, xmin=0, xmax=126,
  xlabel={Kesit boyunca uzaklık [km]}, ylabel={Derinlik [cm]},
  y dir=reverse, ymin=0, ymax=40, grid=none,
  colorbar, colormap/hot, point meta min=20, point meta max=220,
  colorbar style={ylabel={$T$ [$^\circ$C]}}]
\addplot[matrix plot*, point meta=explicit, mesh/cols=120]
  table[x=s,y=z,meta=T] {\dataroot/geo_section.dat};
\end{axis}
\end{tikzpicture}

\caption{Batı--doğu kesiti. Üstte arazi profili, altta aynı kesit boyunca
bağlaşık modelin ürettiği yeraltı sıcaklık alanı.}
\label{fig:geosection}
\end{figure}

% =============================================================================
\section{Enerji ve haberleşme fizibilitesi}

\subsection{Güç bütçesi ve stokastik güneş}

Donanım raporunun bileşen tablosundan yeniden hesaplanan ortalama akım, yüzde
elli görev döngüsünde \iavgFifty~mA, yüzde onda \iavgTen~mA'dir. Güneş
girdisi deterministik bir tepe güneş saati sayısıyla değil, stokastik olarak
modellenmiştir: enleme bağlı dış-atmosfer günlük ışınımı analitik olarak
hesaplanmış, açıklık indisi mevsimsel ortalamaya birinci mertebe
özbağlanımlı bir sapma eklenerek üretilmiş ve üç yıllık bir seri boyunca şarj
durumu izlenmiştir. Kış ve yaz için elde edilen eşdeğer tepe güneş saatleri
\pshWinter\ ve \pshSummer\ saat/gündür.

Çözümlemenin ilk sonucu, tasarım sorusunun yanlış eksende sorulduğunu
göstermiştir. Beş vatlık panel, açık gökyüzü referans koşullarında kış
ortasında bile yükün yaklaşık \energyMargin\ katını üretmektedir; panel gücü
bu nedenle kısıtlayıcı değişken değildir. Kısıtlayıcı olan, kanopi altında
panele ulaşan ışık payıdır. Şekil~\ref{fig:energy}a bu yüzden yük kaybı
olasılığını görev döngüsü ile kanopi geçirgenliği düzleminde verir. Referans
geçirgenlikte (yüzde elli beş) yük kaybı olasılığı yüzde elli görev döngüsünde
\lolpFifty~\%, yüzde onda \lolpTen~\%'dir. Kapalı bir çam örtüsüne karşılık
gelen yüzde on geçirgenlikte ise yüzde elli görev döngüsünde \lolpShade~\%'ye
çıkar. Yüzde bir yük kaybı olasılığının altında kalmak için gereken en düşük
geçirgenlik, yüzde elli görev döngüsünde \canopyCritFifty~\%, kesintisiz
çalışmada \canopyCritFull~\%'dir.

Bunun kurulum açısından karşılığı nettir: enerji fizibilitesi panel seçimine
değil, yüzey düğümünün gölgelenme derecesine bağlıdır ve montaj yönergesi bu
eşiği bir kabul ölçütü olarak içermelidir. Donanım raporunun kendi önerisi olan
olay-tetikli uyarlamalı görev döngüsü, aynı grafikte sakin dönemde soldaki
geniş güvenli bölgede, şüphe hâlinde kısa süreliğine sağda çalışmak anlamına
gelir ve en yoğun gölgelenme koşullarında bile marjı korur.

\begin{figure}[htbp]\centering
\begin{tikzpicture}
\begin{groupplot}[group style={group size=2 by 1, horizontal sep=2.6cm},
  aegis, width=5.2cm, height=5.0cm]
\nextgroupplot[xlabel={Görev döngüsü}, ylabel={Kanopi geçirgenliği},
  title={(a) Yük kaybı olasılığı}, grid=none,
  colorbar, colormap/viridis, point meta min=0, point meta max=0.4]
\addplot[matrix plot*, point meta=explicit, mesh/cols=40]
  table[x=duty,y=canopy,meta=lolp] {\dataroot/energy_lolp.dat};
\nextgroupplot[xlabel={Gün}, ymin=0, ymax=1.06,
  xmin=0, xmax=365, title={(b) Şarj durumu, \%50 görev döngüsü},
  legend style={at={(0.03,0.05)}, anchor=south west}]
\addplot[aegisteal, thick] table[x=day,y=soc_ref] {\dataroot/energy_soc.dat};
\addlegendentry{referans gölgelenme}
\addplot[aegisember, thick] table[x=day,y=soc_shade] {\dataroot/energy_soc.dat};
\addlegendentry{derin gölge (\%4)}
\end{groupplot}
\end{tikzpicture}

\caption{Solda kanopi geçirgenliği--görev döngüsü düzleminde yük kaybı
olasılığı; sağda referans koşulda bir yıllık şarj durumu izi.}
\label{fig:energy}
\end{figure}

\subsection{LoRa bağlantı bütçesi}

Haberleşme fizibilitesi iki katmanlı bir yol kaybı modeliyle değerlendirilmiştir
(Şekil~\ref{fig:link}). Yüz metreye kadar serbest uzay kaybı, ötesinde engebeli
ve örtülü arazi için $n = \pathExp$ üstelli logaritmik uzaklık modeli
kullanılmış, buna Weissberger'in değiştirilmiş üstel sönüm bağıntısıyla bitki
örtüsü sönümü eklenmiştir. Yirmi iki dBm çıkış gücü, iki dBi anten kazançları
ve $-134$~dBm alıcı duyarlılığı ile bağlantı marjının sıfıra düştüğü mesafe
\dMax~metre, on dB marj bırakıldığında \dMarginTen~metredir. Serbest uzay
eğrisiyle arasındaki fark, menzilin neden yalnızca verici gücüyle
hesaplanamayacağını göstermektedir.

Ağ geçidi kapasitesi iki ayrı kısıttan gelir. Menzil kısıtı, yüz metrelik
düğüm aralığında yaklaşık \nodesRange\ düğüme karşılık gelir; hava süresi
kısıtı ise, düğüm başına on dakikada bir yüz milisaniyelik paket ve sekiz
kanalda yüzde yirmi doluluk hedefiyle \nodesAirtime\ düğümdür. Bağlayıcı olan
küçük olanıdır, yani \nodesPerGw\ düğüm. Her iki değer de bir ağ geçidinin tek
başına yüzlerce hektarlık bir parseli taşıyabileceğini gösterir; haberleşme
katmanı kapsama açısından kısıtlayıcı bileşen değildir.

\begin{figure}[htbp]\centering
\begin{tikzpicture}
\begin{axis}[aegis, xlabel={Mesafe [m]}, ylabel={Alınan güç [dBm]},
  xmode=log, xmin=5, xmax=4000, ymin=-165, ymax=-20,
  width=9.6cm, height=5.0cm, legend pos=north east]
\addplot[aegisblue, thick] table[x=d,y=prx_open] {\dataroot/link.dat};
\addlegendentry{açık arazi (serbest uzay)}
\addplot[aegisteal, thick] table[x=d,y=prx_forest] {\dataroot/link.dat};
\addlegendentry{orman örtüsü (Weissberger)}
\addplot[aegisember, dashed, domain=5:4000] {-134};
\addlegendentry{alıcı duyarlılığı $-134$ dBm}
\end{axis}
\end{tikzpicture}

\caption{LoRa bağlantı bütçesi. Bitki örtüsü sönümü ve engebeli arazi yol
kaybı üsteli birlikte, açık arazi menzilini bir mertebe kısaltmaktadır.}
\label{fig:link}
\end{figure}

% =============================================================================
\section{Ekonomik fizibilite}

Ekonomik katman, fiziksel katmandan bağımsız değildir: düğüm yoğunluğu
Bölüm~\ref{sec:lateral}'ten, düğüm maliyeti donanım raporunun malzeme
listesinden, ağ geçidi başına düğüm sayısı bağlantı bütçesinden gelir. Yedi
yıllık bir ufukta, \econRuns\ örneklik bir Monte Carlo tasarımıyla donanım
maliyeti, kâr marjı, yazılım abonelik geliri, müşteri kaybı oranı, iskonto
oranı, büyüme çarpanı, işletme gideri ve ilk yıl kurulum hacmi üçgen
dağılımlarla örneklenmiştir.

Net bugünkü değerin medyanı \npvMed\ milyon TL, onuncu ve doksanıncı
yüzdelikleri \npvPten\ ile \npvPninety\ milyon TL'dir; pozitif çıkma olasılığı
\pNpvPos~\%'dir. İç verim oranının medyanı yüzde \irrMed, onuncu yüzdeliği
yüzde \irrPten'dir. Medyan senaryoda kümülatif iskontolu nakit akışı
\paybackYear{}. yılda pozitife dönmektedir (Şekil~\ref{fig:econ}).

Duyarlılık çözümlemesi (Şekil~\ref{fig:tornado}), sonucun en çok büyüme
çarpanına ve yazılım abonelik gelirine, en az donanım birim maliyetine bağlı
olduğunu göstermektedir. Bu, iş planındaki tekrarlayan gelir vurgusunu
destekler ve donanım maliyetini düşürmeye yönelik erken optimizasyonun
ekonomik getirisinin sınırlı olduğunu söyler.

\begin{figure}[htbp]\centering
\begin{tikzpicture}
\begin{groupplot}[group style={group size=2 by 1, horizontal sep=1.9cm},
  aegis, width=5.4cm, height=5.0cm]
\nextgroupplot[xlabel={NBD [milyon TL]}, ylabel={Frekans}, ybar interval,
  title={(a) Net bugünkü değer}, ymin=0]
\addplot[fill=aegisteal!35, draw=aegisdeep] table[x=npv,y=count] {\dataroot/npv_hist.dat};
\nextgroupplot[xlabel={Yıl}, ylabel={Kümülatif iskontolu nakit [milyon TL]},
  title={(b) Başabaş noktası}]
\addplot[aegisteal, thick, mark=*, mark size=1.6] table[x=year,y=cum] {\dataroot/cashflow.dat};
\addplot[aegismuted, dashed, domain=1:7] {0};
\end{groupplot}
\end{tikzpicture}

\caption{Net bugünkü değerin Monte Carlo dağılımı ve medyan senaryoda
kümülatif iskontolu nakit akışı.}
\label{fig:econ}
\end{figure}

\begin{figure}[htbp]\centering
\begin{tikzpicture}
\begin{axis}[aegis, xbar, y dir=reverse, width=8.4cm, height=5.6cm,
  xlabel={NBD [milyon TL]}, ytick=data,
  yticklabels={donanım maliyeti, kâr marjı, SaaS geliri, kayıp oranı,
               iskonto oranı, büyüme, işletme gideri, ilk yıl kurulum},
  y tick label style={font=\scriptsize}, bar width=6pt, enlarge y limits=0.09,
  legend pos=south east]
\addplot[fill=aegisember!60, draw=aegisember] table[x=low,y=idx] {\dataroot/tornado.dat};
\addlegendentry{p10}
\addplot[fill=aegisteal!60, draw=aegisdeep] table[x=high,y=idx] {\dataroot/tornado.dat};
\addlegendentry{p90}
\end{axis}
\end{tikzpicture}

\caption{Net bugünkü değerin girdi belirsizliklerine duyarlılığı. Her çubuk
çifti, ilgili girdi onuncu ve doksanıncı yüzdeliğine sabitlenip diğerleri
medyanda tutulduğunda elde edilen net bugünkü değeri gösterir.}
\label{fig:tornado}
\end{figure}

% =============================================================================
\section{Bütünleşik değerlendirme}

Tablo~\ref{tab:verdict}, fizibilite ölçütlerini, her biri için önceden
belirlenmiş eşiği ve bu çözümlemenin ürettiği değeri bir arada verir.

\begin{table}[htbp]\centering\small
\caption{Fizibilite ölçütleri ve sayısal sonuçlar.}
\label{tab:verdict}
\begin{tabular}{@{}p{5.1cm} p{3.1cm} p{3.4cm} p{3.3cm}@{}}
\toprule
Ölçüt & Eşik & Hesaplanan & Değerlendirme\\
\midrule
Alev öncesi öncü süre (medyan) & $\ge 15$ dk & \leadMed\ dk &
karşılanıyor\\
Öncü sürenin eşiğin üzerinde kalma olasılığı & $\ge \%80$ &
\pLeadFifteen~\% & karşılanıyor\\
Ayırt edicilik (eğri altı alan) & $\ge 0{,}90$ & \aucFusion &
karşılanıyor\\
$\beta$ eşiğinin tasarım bandı & $\ge 2\times$ genişlik & \betaLo--\betaHi\
[$^\circ$C/10 dk] & karşılanıyor\\
Yanlış alarm oranı ($\tau=1$) & $\le \%1$ & \fprOne~\% & karşılanıyor\\
Sayısal belirsizlik & $\le \%5$ & \gciGrid~\% & karşılanıyor\\
Enerji: kritik kanopi geçirgenliği & $\le \%15$ & \canopyCritFifty~\%
(\%50 görev döngüsü) & montaj ölçütü olarak yazılmalı\\
Haberleşme menzili & $\ge 500$ m & \dMarginTen\ m (10 dB marj) &
karşılanıyor\\
Net bugünkü değerin pozitif olma olasılığı & $\ge \%70$ & \pNpvPos~\% &
karşılanıyor\\
Alev öncesi mekânsal kapsama & --- & $\lambda$'ya bağlı &
\textsc{ölçüm bekliyor}\\
\bottomrule
\end{tabular}
\end{table}

Tablodaki tek açık kalem, mekânsal kapsamadır ve bu bir eksiklik değil, doğru
tanımlanmış bir araştırma sorusudur. Miselyum ağının elektrotonik uzunluk
sabiti ölçülene kadar, sistemin alev öncesi tespit kapsaması bir aralık olarak
kalır. Bu ölçüm, WP-2 iklim odası düzeneğinde ek donanım gerektirmeden
yapılabilir: ağ üzerinde bilinen uzaklıklara yerleştirilen elektrot çiftlerinde
aynı uyarana verilen yanıtın genliği kaydedilir ve denklem~\eqref{eq:cable}
uydurulur. Aynı düzenek, $\gamma$ eşiğinin ve kuruma cephesi hızının
kalibrasyonunu da sağlar.

Bu çözümlemenin doğrudan tasarıma dönüşen üç çıktısı vardır. Birincisi, füzyon
kuralının ikinci koşulu mutlak değer üzerinden yazılmalıdır; bu, yazılımda tek
satırlık bir değişiklik olup altı santimetre derinlikte öncü süreyi
altı santimetrede alarmsızlıktan \leadNom\ dakikaya taşımaktadır. İkincisi, sıcaklık ve direnç probu ile elektrot açıklığı
birlikte ölçeklenerek beş ile sekiz santimetre bandına alınmalı, biyoelektrot
miselyumun yaşadığı derinlikte bırakılmalıdır. Üçüncüsü, görev döngüsü sabit
değil olay-tetikli olmalıdır; enerji fizibilitesi ancak bu koşulda sağlanır.

% =============================================================================
\section{Sınırlar}

Bu çözümlemenin geçerlilik sınırları açıkça belirtilmelidir. Buhar taşınımı
modeli, de Vries çerçevesinin geçerli olduğu kabul edilen sıcaklık bandının
üzerinde, gaz fazı basınç gradyanlarını ve zorlanmış taşınımı ihmal eder;
kaynama cephesinin ötesindeki bölgede model niceliksel değil niteliksel
olarak doğrudur. Toprak parametreleri tek bir dokusal sınıfa aittir ve
mekânsal değişkenlikleri yalnızca ıslaklık indisi üzerinden temsil edilir.
Biyoelektrik kanal, spike istatistiği düzeyinde fenomenolojiktir ve bu
belgede karar kuralının bir parçası olarak değil, üçüncü bir bağımsız kanal
olarak ele alınmıştır; ayırt edicilik hesapları yalnızca iki fiziksel kanal
üzerinden yapılmıştır, dolayısıyla raporlanan eğri altı alanı bir alt sınırdır.
Yüzey sıcaklık senaryoları ölçülmüş yangın verisinden değil, literatürdeki
piroliz sıcaklık aralıklarından kurulmuştur. Ekonomik katmandaki üçgen
dağılımlar uzman görüşüne dayanır ve piyasa verisiyle güncellenmelidir.
Coğrafi katmanda kullanılan duyarlılık indisi, yanmış alan istatistikleriyle
kalibre edilmemiş fiziksel bir birleşimdir; risk haritası olarak değil, saha
seçimi için sıralama aracı olarak yorumlanmalıdır. Aynı katmandaki yerleşim
noktaları, yalnızca büyük merkezleri içeren küçük ölçekli bir veri kümesinden
gelmektedir; yerleşim--orman arayüzü mesafesi bu nedenle üst sınır niteliğinde
olup, saha kararı öncesinde büyük ölçekli bir yerleşim katmanıyla
yenilenmelidir. Yamaç yönelimli anizotropi, coğrafi koşumlarda tam iki boyutlu
tensör çözücüsüyle değil, tensörün düşey bileşeni üzerinden sütun modeline
aktarılarak temsil edilmiştir; yanal akı bileşeninin etkisi
Bölüm~
ef{sec:lateral}'te ayrıca ölçülmüştür ve metre mertebesindedir.

Bu sınırların hiçbiri sonucun işaretini değiştirecek nitelikte değildir. Alev
öncesi imzanın varlığı, ayırt ediciliği ve enerji uygulanabilirliği modelin
geniş bir parametre bandında korunmaktadır; belirsizliğin yoğunlaştığı yer,
sistemin ne kadar erken uyardığı değil, ne kadar geniş bir alanı bu erkenlikle
örtebildiğidir.

% =============================================================================
\section*{Kaynaklar}
\addcontentsline{toc}{section}{Kaynaklar}
\small
\begin{enumerate}\setlength{\itemsep}{0pt}
\item J.~R. Philip, D.~A. de Vries, ``Moisture movement in porous materials
under temperature gradients,'' \emph{Trans. Am. Geophys. Union}, 38(2),
222--232, 1957.
\item M.~Th. van Genuchten, ``A closed-form equation for predicting the
hydraulic conductivity of unsaturated soils,'' \emph{Soil Sci. Soc. Am. J.},
44(5), 892--898, 1980.
\item R.~F. Carsel, R.~S. Parrish, ``Developing joint probability
distributions of soil water retention characteristics,'' \emph{Water Resour.
Res.}, 24(5), 755--769, 1988.
\item J.~Côté, J.-M. Konrad, ``A generalized thermal conductivity model for
soils and construction materials,'' \emph{Can. Geotech. J.}, 42(2), 443--458,
2005.
\item A.~Cass, G.~S. Campbell, T.~L. Jones, ``Enhancement of thermal water
vapor diffusion in soil,'' \emph{Soil Sci. Soc. Am. J.}, 48(1), 25--32, 1984.
\item R.~J. Millington, J.~P. Quirk, ``Permeability of porous solids,''
\emph{Trans. Faraday Soc.}, 57, 1200--1207, 1961.
\item G.~E. Archie, ``The electrical resistivity log as an aid in determining
some reservoir characteristics,'' \emph{Trans. AIME}, 146(1), 54--62, 1942.
\item D.~B. Geselowitz, ``An application of electrocardiographic lead theory
to impedance plethysmography,'' \emph{IEEE Trans. Biomed. Eng.}, BME-18(1),
38--41, 1971.
\item P.~J. Roache, ``Quantification of uncertainty in computational fluid
dynamics,'' \emph{Annu. Rev. Fluid Mech.}, 29, 123--160, 1997.
\item A.~Saltelli ve diğ., ``Variance based sensitivity analysis of model
output,'' \emph{Comput. Phys. Commun.}, 181(2), 259--270, 2010.
\item M.~J.~W. Jansen, ``Analysis of variance designs for model output,''
\emph{Comput. Phys. Commun.}, 117(1--2), 35--43, 1999.
\item R.~D. Barnes, C.~Lehman, D.~Mulla, ``Priority-flood: an optimal
depression-filling and watershed-labeling algorithm,'' \emph{Comput. Geosci.},
62, 117--127, 2014.
\item B.~K.~P. Horn, ``Hill shading and the reflectance map,''
\emph{Proc. IEEE}, 69(1), 14--47, 1981.
\item L.~Kumar, A.~K. Skidmore, E.~Knowles, ``Modelling topographic variation
in solar radiation in a GIS environment,'' \emph{Int. J. Geogr. Inf. Sci.},
11(5), 475--497, 1997.
\item J.~A. Duffie, W.~A. Beckman, \emph{Solar Engineering of Thermal
Processes}, 4.~baskı, Wiley, 2013.
\item M.~Weissberger, \emph{An initial critical summary of models for
predicting the attenuation of radio waves by trees}, ESD-TR-81-101, 1982.
\item R.~C. Rothermel, \emph{A mathematical model for predicting fire spread
in wildland fuels}, USDA Forest Service INT-115, 1972.
\item S.~H. Doerr ve diğ., ``Soil heating during wildfires and prescribed
burns: a global evaluation,'' \emph{Int. J. Wildland Fire}, 34(12), WF25103,
2025.
\item N.~Phillips, A.~Gandia, A.~Adamatzky, ``Electrical response of fungi to
changing moisture content,'' \emph{Fungal Biol. Biotechnol.}, 10(1):8, 2023.
\item A.~Adamatzky, ``Towards fungal computer,'' \emph{Interface Focus},
8(6):20180029, 2018.
\end{enumerate}

\end{document}
